\section{บทที่ 8 โปรแกรมคำนวณทางคณิตศาสตร์สำหรับวิศวกรรมไฟฟ้า}

\textbf{ชื่อบทภาษาอังกฤษ:} Mathematical Computing Software for Electrical Engineering \\
\textbf{รายวิชา:} คณิตศาสตร์วิศวกรรมไฟฟ้า (Electrical Engineering Mathematics) \\
\textbf{รหัสวิชา:} 5571102 \\
\textbf{ระดับผู้เรียน:} ปริญญาตรี สาขาวิชาเทคโนโลยีวิศวกรรมไฟฟ้า \\
\textbf{ช่วงเวลาตามแผนรายวิชา:} สัปดาห์ที่ 14–15 \\
\textbf{จำนวนชั่วโมงที่แนะนำ:} 6 ชั่วโมง (ข้อเสนอเพื่อการออกแบบการสอนสำหรับ 2 สัปดาห์; จำนวนชั่วโมงทางการไม่ได้ระบุในไฟล์โครงสร้างรายวิชา) \\
\textbf{ความยาวเป้าหมาย:} 25–30 หน้า A4 เมื่อจัดหน้าด้วย TH Sarabun New หรือเทียบเท่า ขนาดประมาณ 16 pt \\
\textbf{รูปแบบการเรียน:} บรรยายและฝึกปฏิบัติด้วยโปรแกรมคำนวณ \\

\begin{quote}
\textbf{หมายเหตุเกี่ยวกับบทนี้:} บทที่ 8 เป็นบทบูรณาการและบทปฏิบัติรวบยอดของรายวิชา เน้นการนำหัวข้อจากบทที่ 2–7 ได้แก่ จำนวนเชิงซ้อนและเฟสเซอร์ การวิเคราะห์เวกเตอร์ อนุกรมฟูเรียร์ การแปลงลาปลาซ และการวิเคราะห์วงจร มาตรวจสอบและสำรวจด้วยโปรแกรมคำนวณ ไม่ได้มุ่งสอนการเขียนโปรแกรมเชิงซอฟต์แวร์อย่างลึกซึ้ง
\end{quote}

\par\noindent\rule{\textwidth}{0.4pt}\par

\subsection{1. แนวคิดสำคัญของบท}

งานวิศวกรรมไฟฟ้าสมัยใหม่ต้องอาศัยทั้ง \textbf{แบบจำลองทางคณิตศาสตร์ (Mathematical Model)} และ \textbf{เครื่องมือคำนวณ (Computational Tool)} วิศวกรอาจเริ่มจากกฎของโอห์ม สมการเชิงอนุพันธ์ ฟังก์ชันเชิงซ้อน อนุกรมฟูเรียร์ หรือการแปลงลาปลาซ แล้วใช้โปรแกรมช่วยจัดการการคำนวณที่มีจำนวนข้อมูลมาก มีขั้นตอนซ้ำ หรือมีสมการที่ซับซ้อนกว่าที่เหมาะกับการคำนวณด้วยมือ

อย่างไรก็ตาม โปรแกรมคำนวณไม่ใช่สิ่งทดแทนความเข้าใจทางคณิตศาสตร์ ผลลัพธ์จากโปรแกรมจะมีคุณค่าก็ต่อเมื่อผู้ใช้สามารถ

\begin{enumerate}
    \item สร้างแบบจำลองหรือสมการได้ถูกต้อง
    \item เลือกเครื่องมือและวิธีคำนวณที่เหมาะสม
    \item กำหนดหน่วย สมมติฐาน และช่วงข้อมูลได้ถูกต้อง
    \item อ่านผลลัพธ์ กราฟ และสเปกตรัมได้อย่างมีความหมาย
    \item ตรวจสอบผลลัพธ์ด้วยการประมาณค่า การคำนวณด้วยมือ หรือหลักทางกายภาพ
    \item รายงานวิธีคำนวณและผลลัพธ์ให้ผู้อื่นตรวจสอบซ้ำได้
\end{enumerate}

บทนี้จึงใช้แนวคิด \textbf{Model → Compute → Visualize → Verify → Interpret → Report} เป็นแกนกลางของการเรียนรู้

% [รูปที่ 1: วงจรการแก้ปัญหาทางวิศวกรรมด้วยเครื่องมือคำนวณ]
% ประเภทภาพ: Flowchart / Infographic
% ตำแหน่งที่แนะนำ: หลังหัวข้อ 1 แนวคิดสำคัญของบท
% วัตถุประสงค์การเรียนรู้: ช่วยให้ผู้เรียนเห็นว่าโปรแกรมเป็นส่วนหนึ่งของกระบวนการแก้ปัญหา ไม่ใช่จุดเริ่มต้นหรือคำตอบสุดท้าย
% คำอธิบายภาพ: แผนภาพวงจร 6 ขั้น ได้แก่ Physical Problem → Mathematical Model → Computational Method → Visualization → Verification → Engineering Interpretation & Report และมีลูกศรวนกลับจาก Verification ไปยัง Model เมื่อผลไม่สมเหตุสมผล
% องค์ประกอบที่ต้องมี: ปัญหาจริง, สมการ, เครื่องคิดเลข/Python/MATLAB/Octave/GeoGebra, กราฟ, เครื่องหมายตรวจสอบ, รายงาน
% ข้อความหรือป้ายกำกับ: “ปัญหาทางกายภาพ”, “แบบจำลอง”, “คำนวณ”, “แสดงผล”, “ตรวจสอบ”, “แปลผลและรายงาน”
% ภาษาของป้ายกำกับ: ไทยและอังกฤษ
% สัดส่วนภาพ: 16:9
% Image Generation Prompt: Clean engineering education flowchart showing a six-stage computational problem-solving cycle for electrical engineering: Physical Problem, Mathematical Model, Computational Method, Visualization, Verification, Engineering Interpretation and Report. Include icons for an electrical circuit, equations, scientific calculator, Python/MATLAB/Octave, plots, check mark, and technical report. Add a feedback arrow from Verification back to Mathematical Model. Bilingual Thai-English labels, white background, precise vector textbook style, 16:9.
% Alt Text: แผนภาพกระบวนการแก้ปัญหาทางวิศวกรรมจากปัญหาจริงไปสู่แบบจำลอง การคำนวณ การสร้างกราฟ การตรวจสอบ และการรายงานผล

\begin{figure}[htbp]
    \centering
    \includegraphics[width=0.8\textwidth]{images/figure-01.png}
    \caption{วงจรการแก้ปัญหาทางวิศวกรรมด้วยเครื่องมือคำนวณ}
\end{figure}

\par\noindent\rule{\textwidth}{0.4pt}\par

\subsection{2. ความรู้พื้นฐานก่อนเรียน}

ก่อนเรียนบทนี้ นักศึกษาควรมีความรู้และทักษะดังต่อไปนี้

\begin{itemize}
    \item พีชคณิต ฟังก์ชัน กราฟ สมการ และหน่วยทางวิศวกรรม
    \item จำนวนเชิงซ้อน รูปสี่เหลี่ยม รูปเชิงขั้ว และเฟสเซอร์
    \item อิมพีแดนซ์ของตัวต้านทาน ตัวเหนี่ยวนำ และตัวเก็บประจุ
    \item เวกเตอร์ ดอตโปรดักต์ ครอสโปรดักต์ เกรเดียนต์ ไดเวอร์เจนซ์ และเคิร์ลในระดับพื้นฐาน
    \item ฟังก์ชันเป็นคาบ อนุกรมฟูเรียร์ ฮาร์มอนิก และสเปกตรัม
    \item นิยามและสมบัติพื้นฐานของการแปลงลาปลาซ
    \item วงจร RC, RL, RLC ฟังก์ชันถ่ายโอน โพล และซีโรในระดับเบื้องต้น
    \item การอ่านสูตรและการแทนค่าด้วยหน่วย SI
    \item การใช้งานคอมพิวเตอร์เบื้องต้น เช่น การสร้างไฟล์ การใช้โฟลเดอร์ และการบันทึกผลลัพธ์
\end{itemize}

\textbf{แนวคิดสำคัญก่อนเริ่ม:} นักศึกษาควรสามารถคำนวณตัวอย่างขนาดเล็กด้วยมือได้ก่อน แล้วจึงใช้โปรแกรมขยายไปสู่กรณีที่มีข้อมูลจำนวนมากหรือจำเป็นต้องสร้างกราฟ

\par\noindent\rule{\textwidth}{0.4pt}\par

\subsection{3. คำสำคัญประจำบท}

\begin{table}[htbp]
    \centering
    \begin{tabular}{l l l}
        \hline
        \textbf{คำภาษาไทย} & \textbf{English} & \textbf{ความหมายโดยย่อ} \\
        \hline
        โปรแกรมคำนวณทางคณิตศาสตร์ & Mathematical Computing Software & โปรแกรมที่ช่วยคำนวณ วิเคราะห์ และแสดงผลข้อมูลทางคณิตศาสตร์ \\
        การคำนวณเชิงตัวเลข & Numerical Computation & การหาคำตอบในรูปตัวเลขโดยใช้อัลกอริทึม \\
        การคำนวณเชิงสัญลักษณ์ & Symbolic Computation & การจัดรูป หาอนุพันธ์ ปริพันธ์ หรือทรานส์ฟอร์มโดยคงตัวแปรเชิงสัญลักษณ์ \\
        จำนวนเชิงซ้อน & Complex Number & จำนวนรูป $a+jb$ ที่มีส่วนจริงและส่วนจินตภาพ \\
        เฟสเซอร์ & Phasor & ตัวแทนเชิงซ้อนของสัญญาณไซน์ที่ความถี่คงที่ \\
        อาร์เรย์ & Array & โครงสร้างข้อมูลตัวเลขหนึ่งมิติหรือหลายมิติ \\
        เวกเตอร์ไรซ์ & Vectorization & การดำเนินการกับข้อมูลเป็นชุดแทนการคำนวณทีละสมาชิก \\
        การสร้างกราฟ & Visualization / Plotting & การแสดงข้อมูลหรือฟังก์ชันในรูปกราฟเพื่อช่วยตีความ \\
        FFT & Fast Fourier Transform & อัลกอริทึมที่คำนวณ DFT ได้อย่างมีประสิทธิภาพ \\
        สเปกตรัม & Spectrum & การแสดงองค์ประกอบของสัญญาณตามความถี่ \\
        การแปลงลาปลาซ & Laplace Transform & การแปลงฟังก์ชันจากโดเมนเวลาไปสู่โดเมน $s$ \\
        ฟังก์ชันถ่ายโอน & Transfer Function & อัตราส่วนเอาต์พุตต่ออินพุตในโดเมนลาปลาซภายใต้เงื่อนไขเริ่มต้นเป็นศูนย์ \\
        การตรวจสอบผล & Verification & การตรวจว่าผลคำนวณสอดคล้องกับทฤษฎี หน่วย และพฤติกรรมทางกายภาพ \\
        ความสามารถทำซ้ำ & Reproducibility & ความสามารถในการทำตามขั้นตอนเดิมแล้วได้ผลลัพธ์สอดคล้องกัน \\
        \hline
    \end{tabular}
\end{table}

\par\noindent\rule{\textwidth}{0.4pt}\par

\subsection{4. ผลลัพธ์การเรียนรู้ประจำบท}

เมื่อเรียนจบบทนี้ นักศึกษาสามารถ...

\begin{enumerate}
    \item \textbf{เลือกใช้} เครื่องคิดเลขวิศวกรรม Python, MATLAB, GNU Octave หรือ GeoGebra ให้เหมาะกับลักษณะปัญหาทางวิศวกรรมไฟฟ้าได้
    \item \textbf{คำนวณและแสดงผล} จำนวนเชิงซ้อน เฟสเซอร์ และเวกเตอร์ด้วยโปรแกรมได้อย่างถูกต้อง พร้อมกำหนดหน่วยและรูปแบบผลลัพธ์ที่เหมาะสม
    \item \textbf{สร้างและวิเคราะห์} รูปคลื่น อนุกรมฟูเรียร์ และสเปกตรัมด้วยวิธีเชิงตัวเลขได้ พร้อมอธิบายผลของอัตราการสุ่มตัวอย่างและการกำหนดแกนความถี่
    \item \textbf{ประยุกต์ใช้} การคำนวณเชิงสัญลักษณ์และเชิงตัวเลขกับการแปลงลาปลาซ ฟังก์ชันถ่ายโอน และการตอบสนองของวงจร RC/RL/RLC ได้
    \item \textbf{ตรวจสอบและเปรียบเทียบ} ผลลัพธ์จากโปรแกรมกับการคำนวณด้วยมือหรือข้อจำกัดทางกายภาพ และระบุความคลาดเคลื่อนที่สำคัญได้
    \item \textbf{จัดทำรายงาน} ผลการคำนวณทางวิศวกรรมไฟฟ้าที่ระบุสมการ พารามิเตอร์ หน่วย วิธีคำนวณ กราฟ และข้อสรุปอย่างเป็นระบบได้
\end{enumerate}

\subsubsection{4.1 ตารางเชื่อมโยง OBE}

\begin{table}[htbp]
    \centering
    \begin{tabular}{l l l l l}
        \hline
        \textbf{ผลลัพธ์การเรียนรู้ประจำบท} & \textbf{เนื้อหาที่เกี่ยวข้อง} & \textbf{กิจกรรมการเรียนรู้} & \textbf{วิธีประเมิน} & \textbf{หลักฐานการเรียนรู้} \\
        \hline
        CLO1 เลือกใช้เครื่องมือ & 8.1–8.6 & วิเคราะห์กรณีศึกษาและเลือกเครื่องมือ & คำถามสั้น/ใบงาน & ตารางเหตุผลการเลือกเครื่องมือ \\
        CLO2 คำนวณจำนวนเชิงซ้อนและเวกเตอร์ & 8.7–8.8 & ทดลองคำนวณด้วยมือและโปรแกรม & แบบฝึกหัดปฏิบัติ & โค้ด/ไฟล์คำนวณและคำตอบที่ตรวจหน่วยแล้ว \\
        CLO3 วิเคราะห์ฟูเรียร์และสเปกตรัม & 8.9 & สร้างสัญญาณและ FFT & งานปฏิบัติพร้อมกราฟ & กราฟโดเมนเวลาและสเปกตรัมพร้อมคำอธิบาย \\
        CLO4 ประยุกต์ลาปลาซกับวงจร & 8.10 & Symbolic Laplace + numerical response & แบบฝึกหัด/ชิ้นงาน & สมการ ฟังก์ชันถ่ายโอน โค้ด และกราฟตอบสนอง \\
        CLO5 ตรวจสอบผล & ทุกหัวข้อ & เปรียบเทียบมือ–โปรแกรมและ error analysis & Rubric การตรวจสอบ & ตารางเปรียบเทียบ ค่าคลาดเคลื่อน และการอภิปราย \\
        CLO6 จัดทำรายงาน & 8.11–8.12 & Mini computational report & ประเมินรายงาน & รายงาน 2–4 หน้า พร้อมไฟล์ต้นฉบับ \\
        \hline
    \end{tabular}
\end{table}

\par\noindent\rule{\textwidth}{0.4pt}\par

\subsection{5. แผนผังเนื้อหาของบท}

ลำดับการเรียนรู้ของบทนี้ควรเป็นดังนี้

\begin{enumerate}
    \item \textbf{เข้าใจบทบาทของเครื่องมือ} และข้อจำกัดของการคำนวณด้วยโปรแกรม
    \item \textbf{เลือกเครื่องมือให้ตรงงาน} เริ่มจากเครื่องคิดเลข → Python → MATLAB/Octave → GeoGebra
    \item \textbf{ทบทวนการคำนวณพื้นฐาน} จำนวนเชิงซ้อน เฟสเซอร์ และเวกเตอร์
    \item \textbf{เพิ่มมิติของข้อมูล} จากค่าตัวเลขเดี่ยวไปสู่สัญญาณและอาร์เรย์
    \item \textbf{วิเคราะห์ความถี่} ด้วย Fourier series และ FFT
    \item \textbf{วิเคราะห์ระบบในโดเมน $s$} ด้วย Laplace transform และฟังก์ชันถ่ายโอน
    \item \textbf{ตรวจสอบผลลัพธ์} โดยเปรียบเทียบเชิงสัญลักษณ์ เชิงตัวเลข และเหตุผลทางกายภาพ
    \item \textbf{จัดทำรายงาน} ที่ผู้อื่นสามารถตรวจสอบและทำซ้ำได้
\end{enumerate}

% [รูปที่ 2: แผนผังการเลือกเครื่องมือคำนวณตามลักษณะปัญหา]
% ประเภทภาพ: Decision Flowchart
% ตำแหน่งที่แนะนำ: หลังหัวข้อ 5
% วัตถุประสงค์การเรียนรู้: ช่วยเลือกเครื่องมือจากชนิดงาน ไม่ใช่เลือกจากความคุ้นเคยเพียงอย่างเดียว
% คำอธิบายภาพ: เริ่มจากคำถาม “โจทย์ต้องการอะไร?” แตกแขนงเป็น ค่าตัวเลขรวดเร็ว → Scientific Calculator, อาร์เรย์/ข้อมูล/FFT → Python NumPy/SciPy, symbolic → SymPy หรือ MATLAB Symbolic, matrix/engineering workflow → MATLAB/Octave, visualization/interactive concept → GeoGebra
% องค์ประกอบที่ต้องมี: decision diamonds, calculator, Python, MATLAB/Octave, GeoGebra, graph, matrix, symbolic expression
% ข้อความหรือป้ายกำกับ: Numerical, Symbolic, Array/Signal, Interactive Graph, Quick Check
% ภาษาของป้ายกำกับ: ไทยและอังกฤษ
% สัดส่วนภาพ: 16:9
% Image Generation Prompt: Educational decision flowchart for selecting mathematical computing tools in electrical engineering. Start with “What does the problem require?” Branch to Quick Numerical Check → Scientific Calculator; Arrays, Signals and FFT → Python NumPy/SciPy; Symbolic Algebra and Laplace → SymPy or MATLAB Symbolic; Matrix-Based Engineering Workflow → MATLAB or GNU Octave; Interactive Function Exploration → GeoGebra. Clean bilingual Thai-English labels, white background, vector textbook style, 16:9.
% Alt Text: แผนผังตัดสินใจเลือกเครื่องมือคำนวณตามความต้องการของโจทย์

\begin{figure}[htbp]
    \centering
    \includegraphics[width=0.8\textwidth]{images/figure-02.png}
    \caption{แผนผังการเลือกเครื่องมือคำนวณตามลักษณะปัญหา}
\end{figure}

\par\noindent\rule{\textwidth}{0.4pt}\par

\subsection{6. บทนำ}

ในบทก่อนหน้า นักศึกษาใช้คณิตศาสตร์แก้ปัญหาทีละโจทย์ เช่น หาอิมพีแดนซ์ของวงจร AC คำนวณเกรเดียนต์ของศักย์ไฟฟ้า หาอนุกรมฟูเรียร์ของรูปคลื่น หรือใช้การแปลงลาปลาซหาการตอบสนองของวงจร RC เมื่อจำนวนตัวแปรเพิ่มขึ้นหรือจำเป็นต้องสร้างกราฟหลายจุด การคำนวณด้วยมือจะใช้เวลามากและเกิดความผิดพลาดได้ง่าย

โปรแกรมคำนวณช่วยแก้ข้อจำกัดนี้ได้สามรูปแบบหลัก

\begin{itemize}
    \item \textbf{คำนวณเชิงตัวเลข:} เช่น คำนวณอิมพีแดนซ์ 100 ค่าในช่วงความถี่ 10 Hz–10 kHz
    \item \textbf{คำนวณเชิงสัญลักษณ์:} เช่น หา $\mathcal{L}\{e^{-at}\sin \omega t\}$ โดยเก็บ $a$ และ $\omega$ เป็นตัวแปร
    \item \textbf{แสดงผลและสำรวจ:} เช่น ปรับจำนวนฮาร์มอนิก $N$ แล้วสังเกตการเข้าใกล้รูปคลื่นสี่เหลี่ยม
\end{itemize}

ความเสี่ยงสำคัญคือโปรแกรมสามารถให้ “คำตอบที่คำนวณตามคำสั่งได้ถูกต้อง” แม้ว่าผู้ใช้จะป้อนสมการ หน่วย หรือพารามิเตอร์ผิด ดังนั้นบทนี้จึงเน้น \textbf{การตรวจสอบสามชั้น} ได้แก่

\begin{enumerate}
    \item \textbf{ตรวจสมการ:} แบบจำลองสอดคล้องกับระบบจริงหรือไม่
    \item \textbf{ตรวจตัวเลขและหน่วย:} ขนาดและหน่วยสมเหตุสมผลหรือไม่
    \item \textbf{ตรวจพฤติกรรม:} กราฟและแนวโน้มเป็นไปตามหลักทางไฟฟ้าหรือไม่
\end{enumerate}

ตัวอย่างเช่น วงจร RC low-pass เมื่อ $\omega \to 0$ ควรมี $|H(j\omega)|\to 1$ และเมื่อ $\omega\to\infty$ ควรมี $|H(j\omega)|\to 0$ หากกราฟจากโปรแกรมขัดกับพฤติกรรมนี้ ต้องตรวจแบบจำลองหรือคำสั่งก่อนยอมรับผลลัพธ์

\par\noindent\rule{\textwidth}{0.4pt}\par

\section{7. เนื้อหาหลัก}

\subsection{7.1 บทบาทของโปรแกรมคำนวณในงานวิศวกรรมไฟฟ้า}

\subsubsection{7.1.1 งานที่โปรแกรมคำนวณช่วยได้}

โปรแกรมคำนวณมีบทบาทสำคัญในงานวิศวกรรมไฟฟ้า เช่น

\begin{itemize}
    \item คำนวณจำนวนเชิงซ้อนและเฟสเซอร์
    \item แก้ระบบสมการเชิงเส้นในวงจรหลายโหนด
    \item คำนวณเวกเตอร์และสนาม
    \item สร้างรูปคลื่นและวิเคราะห์สเปกตรัม
    \item คำนวณ FFT จากข้อมูลวัด
    \item หาปริพันธ์ อนุพันธ์ และทรานส์ฟอร์มเชิงสัญลักษณ์
    \item แก้สมการเชิงอนุพันธ์เชิงตัวเลข
    \item คำนวณและแสดงการตอบสนองของระบบ LTI
    \item สร้างกราฟ Bode, step response หรือเส้นสนามตามความเหมาะสม
    \item สร้างตารางผลลัพธ์และเตรียมข้อมูลสำหรับรายงาน
\end{itemize}

\subsubsection{7.1.2 Numerical กับ Symbolic ต่างกันอย่างไร}

ให้พิจารณาสมการ

\[
Z_L = j\omega L
\]

หากกำหนด $f=50\ \text{Hz}$ และ $L=50\ \text{mH}$ แล้วคำนวณค่า $Z_L$ เป็นจำนวนเชิงซ้อน นั่นคือ \textbf{การคำนวณเชิงตัวเลข}

แต่ถ้าต้องการเก็บ $\omega$ และ $L$ ไว้เป็นตัวแปรเพื่อจัดรูป เช่น หา

\[
\frac{d}{d\omega}|Z_L|
\]

หรือทำ Laplace transform โดยไม่แทนค่าพารามิเตอร์ นั่นคือ \textbf{การคำนวณเชิงสัญลักษณ์}

\begin{table}[htbp]
    \centering
    \begin{tabular}{l l l}
        \hline
        \textbf{ประเด็น} & \textbf{Numerical} & \textbf{Symbolic} \\
        \hline
        ผลลัพธ์ & ตัวเลขประมาณ & นิพจน์/สูตร \\
        เหมาะกับ & ข้อมูลจำนวนมาก การจำลอง กราฟ & จัดรูปสมการ พิสูจน์ หา transform \\
        ตัวอย่างเครื่องมือ & NumPy, SciPy, MATLAB/Octave & SymPy, MATLAB Symbolic, GeoGebra CAS \\
        จุดแข็ง & เร็วและเหมาะกับข้อมูลจริง & เห็นโครงสร้างทางคณิตศาสตร์ \\
        ข้อควรระวัง & rounding, sampling, conditioning & สมมติฐานของตัวแปร รูปคำตอบอาจซับซ้อน \\
        \hline
    \end{tabular}
\end{table}

\subsubsection{7.1.3 หลัก “คำนวณด้วยมือก่อนขยายด้วยโปรแกรม”}

โจทย์เพื่อการเรียนรู้ควรเริ่มจากกรณีขนาดเล็กที่นักศึกษาคำนวณด้วยมือได้ เช่น

\begin{enumerate}
    \item คำนวณอิมพีแดนซ์ที่ความถี่เดียวด้วยมือ
    \item เขียนโปรแกรมคำนวณอิมพีแดนซ์ในช่วง 100 ความถี่
    \item สร้างกราฟ magnitude และ phase
    \item ตรวจว่าค่าที่ความถี่ที่คำนวณด้วยมือตรงกับโปรแกรม
\end{enumerate}

วิธีนี้ลดการ “เชื่อโปรแกรมโดยไม่ตรวจสอบ” และสร้างทักษะวิศวกรรมที่สำคัญกว่าการจำคำสั่ง

\par\noindent\rule{\textwidth}{0.4pt}\par

\subsection{7.2 การใช้เครื่องคิดเลขวิศวกรรมอย่างมีประสิทธิภาพ}

เครื่องคิดเลขวิศวกรรมยังเป็นเครื่องมือที่มีประโยชน์มากสำหรับการตรวจคำตอบอย่างรวดเร็ว แม้จะมี Python หรือ MATLAB ก็ตาม

\subsubsection{7.2.1 โหมดที่ควรรู้}

\begin{itemize}
    \item DEG / RAD สำหรับมุม
    \item Complex mode สำหรับ $a+jb$
    \item Rectangular ↔ Polar conversion
    \item Scientific notation
    \item Matrix / simultaneous equations หากรุ่นรองรับ
    \item Statistics สำหรับค่าเฉลี่ยและส่วนเบี่ยงเบน
\end{itemize}

\subsubsection{7.2.2 ความผิดพลาดที่พบบ่อย}

\begin{enumerate}
    \item ใช้ DEG แทน RAD ในสูตรที่โปรแกรมรับเรเดียน
    \item ป้อน $50\text{ mH}$ เป็น 50 แทน $50\times10^{-3}$ H
    \item ลืมวงเล็บ เช่น $1/(2\pi fC)$
    \item ใช้ค่าพีแบบปัดเศษมากเกินไปในหลายขั้น
    \item ปัดเลขกลางทางก่อนจบการคำนวณ
\end{enumerate}

\subsubsection{ตัวอย่างที่ 1 การตรวจสอบวงจร RLC อนุกรมด้วยเครื่องคิดเลข}

\textbf{ตัวอย่างเพิ่มเติมเพื่อการเรียนรู้}

ให้วงจร RLC อนุกรมมี

\begin{itemize}
    \item $R=20\ \Omega$
    \item $L=50\ \text{mH}$
    \item $C=100\ \mu\text{F}$
    \item $f=50\ \text{Hz}$
    \item $V=230\ \text{V}_{\text{rms}}$
\end{itemize}

ต้องการหา $Z$ และกระแสเฟสเซอร์ $I$

\textbf{1) วิเคราะห์ระบบ} \\
อิมพีแดนซ์รวมของวงจรอนุกรมคือ

\[
Z = R + j\left(\omega L-\frac{1}{\omega C}\right)
\]

\textbf{2) กำหนดตัวแปร}

\[
\omega=2\pi f=2\pi(50)=314.159\ \text{rad/s}
\]

\textbf{3) คำนวณรีแอกแตนซ์}

\[
X_L=\omega L=314.159(0.05)=15.708\ \Omega
\]

\[
X_C=\frac{1}{\omega C}=\frac{1}{314.159(100\times10^{-6})}=31.831\ \Omega
\]

ดังนั้น

\[
Z=20-j16.123\ \Omega
\]

\textbf{4) แปลงเป็นรูปเชิงขั้ว}

\[
|Z|=\sqrt{20^2+(-16.123)^2}\approx25.690\ \Omega
\]

\[
\angle Z=\tan^{-1}\left(\frac{-16.123}{20}\right)\approx-38.874^\circ
\]

จึงได้

\[
Z\approx25.690\angle -38.874^\circ\ \Omega
\]

\textbf{5) หากกำหนดแรงดันเป็นเฟสอ้างอิง}

\[
V=230\angle0^\circ\ \text{V}
\]

\[
I=\frac{V}{Z}=\frac{230\angle0^\circ}{25.690\angle-38.874^\circ}
\approx8.953\angle38.874^\circ\ \text{A}
\]

ในรูปสี่เหลี่ยม

\[
I\approx6.970+j5.619\ \text{A}
\]

\textbf{6) ตรวจความสมเหตุสมผล} \\
วงจรมี $X_C>X_L$ จึงมีลักษณะเชิงเก็บประจุ กระแสควรนำแรงดัน ซึ่งสอดคล้องกับมุมกระแส $+38.874^\circ$

\begin{quote}
\textbf{ข้อสังเกต:} ตัวอย่างนี้จะถูกใช้ซ้ำในหัวข้อ Python เพื่อยืนยันว่าผลจากโปรแกรมตรงกับการคำนวณด้วยมือ
\end{quote}

\par\noindent\rule{\textwidth}{0.4pt}\par

\subsection{7.3 การใช้ Python สำหรับคำนวณทางวิศวกรรมไฟฟ้า}

Python เหมาะกับการเรียนรู้เพราะสามารถเริ่มจากการคำนวณง่าย ๆ แล้วขยายไปสู่อาร์เรย์ กราฟ สัญญาณ และ symbolic computation ได้ในภาษาเดียว โดยในบทนี้ใช้แนวคิดมากกว่าการผูกติดกับ IDE ใด IDE หนึ่ง

\subsubsection{7.3.1 โครงสร้างพื้นฐานที่จำเป็น}

สิ่งที่นักศึกษาควรรู้เพียงพอสำหรับบทนี้ ได้แก่

\begin{itemize}
    \item ตัวแปรและชนิดข้อมูลตัวเลข
    \item complex number เช่น \texttt{3+4j}
    \item list และ array
    \item การเรียกฟังก์ชัน
    \item การนำเข้าไลบรารีด้วย \texttt{import}
    \item การสร้างกราฟ
    \item การพิมพ์ผลลัพธ์ด้วยรูปแบบที่ควบคุมจำนวนทศนิยม
\end{itemize}

\subsubsection{7.3.2 จำนวนเชิงซ้อนใน Python}

Python ใช้ \texttt{j} แทนหน่วยจินตภาพ เช่น

\begin{verbatim}
z = 3 + 4j
print(z.real)
print(z.imag)
print(abs(z))
\end{verbatim}

สำหรับฟังก์ชันที่ทำงานกับจำนวนเชิงซ้อน สามารถใช้ \texttt{cmath} เช่น phase และการแปลง polar/rectangular

\begin{verbatim}
import cmath
import math

z = 3 + 4j
r, phi = cmath.polar(z)
print(r)
print(math.degrees(phi))
\end{verbatim}

\subsubsection{ตัวอย่างที่ 2 ตรวจสอบวงจร RLC ของตัวอย่างที่ 1 ด้วย Python}

\textbf{ตัวอย่างเพิ่มเติมเพื่อการเรียนรู้}

\begin{verbatim}
import math
import cmath

R = 20.0
L = 50e-3
C = 100e-6
f = 50.0
V = 230.0

w = 2 * math.pi * f
Z = R + 1j * (w*L - 1/(w*C))
I = V / Z

print(f"Z = {Z:.4f} ohm")
print(f"|Z| = {abs(Z):.4f} ohm")
print(f"angle Z = {math.degrees(cmath.phase(Z)):.4f} deg")
print(f"I = {I:.4f} A")
print(f"|I| = {abs(I):.4f} A")
print(f"angle I = {math.degrees(cmath.phase(I)):.4f} deg")
\end{verbatim}

ผลที่ควรได้ใกล้เคียง

\begin{verbatim}
Z = 20.0000-16.1230j ohm
|Z| = 25.6895 ohm
angle Z = -38.8741 deg
I = 6.9702+5.6190j A
|I| = 8.9531 A
angle I = 38.8741 deg
\end{verbatim}

\textbf{การตรวจสอบ}

\begin{itemize}
    \item ขนาดและมุมตรงกับการคำนวณด้วยมือ
    \item กระแสนำแรงดันเพราะวงจรเป็น capacitive
    \item หน่วยต้องกำหนดโดยผู้ใช้ โปรแกรมไม่ทราบโดยอัตโนมัติว่า \texttt{L=50e-3} หมายถึงเฮนรี
\end{itemize}

\subsubsection{7.3.3 จากค่าจุดเดียวไปสู่ช่วงความถี่}

จุดแข็งของโปรแกรมคือการคำนวณแบบอาร์เรย์ เช่น ต้องการดู $|Z|$ ตั้งแต่ 10 Hz ถึง 1 kHz สามารถสร้างช่วงความถี่แล้วคำนวณทุกจุดพร้อมกันด้วย NumPy

\begin{verbatim}
import numpy as np

f = np.logspace(1, 3, 300)       # 10 Hz to 1000 Hz
w = 2*np.pi*f
Z = R + 1j*(w*L - 1/(w*C))
Zmag = np.abs(Z)
\end{verbatim}

จากนั้นจึงนำ \texttt{f} และ \texttt{Zmag} ไปสร้างกราฟได้

\par\noindent\rule{\textwidth}{0.4pt}\par

\subsection{7.4 การใช้ไลบรารี NumPy, SciPy, Matplotlib และ SymPy}

ไลบรารีแต่ละตัวมีหน้าที่ต่างกัน การเลือกใช้ควรสัมพันธ์กับชนิดของงาน

\begin{table}[htbp]
    \centering
    \begin{tabular}{l l l}
        \hline
        \textbf{ไลบรารี} & \textbf{บทบาทหลักในบทนี้} & \textbf{ตัวอย่างงาน} \\
        \hline
        NumPy & อาร์เรย์ คำนวณเชิงตัวเลข FFT & รูปคลื่น อิมพีแดนซ์ตามความถี่ สเปกตรัม \\
        SciPy & อัลกอริทึมวิทยาศาสตร์และ signal/system & step response, frequency response \\
        Matplotlib & กราฟและการแสดงผล & waveform, spectrum, response curve \\
        SymPy & คณิตศาสตร์เชิงสัญลักษณ์ & simplify, derivative, integral, Laplace transform \\
        \hline
    \end{tabular}
\end{table}

\subsubsection{7.4.1 NumPy: อาร์เรย์และเวกเตอร์ไรซ์}

ตัวอย่างสร้างสัญญาณไซน์ 50 Hz เป็นเวลา 0.1 s ที่อัตราสุ่มตัวอย่าง 5 kHz

\begin{verbatim}
import numpy as np

fs = 5000.0
T = 0.1
t = np.arange(0, T, 1/fs)
v = 311*np.sin(2*np.pi*50*t)
\end{verbatim}

จำนวนตัวอย่างโดยประมาณคือ

\[
N=f_sT=5000(0.1)=500
\]

ความสัมพันธ์นี้สำคัญเพราะ FFT จะทำงานบนชุดข้อมูลที่สุ่มตัวอย่างแล้ว ไม่ใช่ฟังก์ชันต่อเนื่องโดยตรง

\subsubsection{7.4.2 Matplotlib: กราฟที่อ่านความหมายได้}

\begin{verbatim}
import matplotlib.pyplot as plt

plt.plot(t, v)
plt.xlabel("Time (s)")
plt.ylabel("Voltage (V)")
plt.title("50-Hz Sinusoidal Voltage")
plt.grid(True)
plt.show()
\end{verbatim}

กราฟวิศวกรรมที่ดีควรมีอย่างน้อย

\begin{itemize}
    \item ชื่อกราฟ
    \item ชื่อแกน
    \item หน่วย
    \item legend เมื่อมีหลายเส้น
    \item ช่วงแกนที่เหมาะสม
    \item ข้อความอธิบายพารามิเตอร์สำคัญในรายงาน
\end{itemize}

% [รูปที่ 3: โครงสร้างข้อมูลจากสมการสู่กราฟใน Python]
% ประเภทภาพ: Block Diagram
% ตำแหน่งที่แนะนำ: หัวข้อ 7.4
% วัตถุประสงค์การเรียนรู้: แสดงความสัมพันธ์ระหว่างพารามิเตอร์ อาร์เรย์เวลา สมการสัญญาณ และกราฟ
% คำอธิบายภาพ: แสดง blocks fs, T → t = np.arange(...) → v(t) → plt.plot(t,v) พร้อมตัวอย่าง 50 Hz sinusoid
% องค์ประกอบที่ต้องมี: sampling frequency, time array, signal array, plot
% ภาษาของป้ายกำกับ: ไทยและอังกฤษ
% สัดส่วนภาพ: 16:9
% Image Generation Prompt: Clean engineering block diagram showing the Python numerical workflow for a sampled electrical waveform: sampling frequency fs and duration T, generate time array t, evaluate v(t)=Vm sin(2πft+phi), store voltage array, then plot time versus voltage with Matplotlib. Include a small sinusoidal output plot. Bilingual Thai-English labels, white background, vector textbook style, 16:9.
% Alt Text: แผนภาพการสร้างอาร์เรย์เวลา คำนวณสัญญาณ และสร้างกราฟด้วย Python

\begin{figure}[htbp]
    \centering
    \includegraphics[width=0.8\textwidth]{images/figure-03.png}
    \caption{โครงสร้างข้อมูลจากสมการสู่กราฟใน Python}
\end{figure}

\subsubsection{7.4.3 SciPy: ระบบเชิงเส้นและการตอบสนอง}

สำหรับระบบ LTI ต่อเนื่อง ถ้าฟังก์ชันถ่ายโอนเขียนเป็น

\[
H(s)=\frac{b_0s^m+\cdots+b_m}{a_0s^n+\cdots+a_n}
\]

สามารถเก็บสัมประสิทธิ์ numerator และ denominator แล้วใช้เครื่องมือใน \texttt{scipy.signal} เพื่อหาการตอบสนองแบบ step หรือ frequency response

ตัวอย่าง RC low-pass

\[
H(s)=\frac{1}{RCs+1}
\]

หาก $R=1\ \text{k}\Omega$ และ $C=100\ \mu\text{F}$ จะได้ $RC=0.1$ s

\begin{verbatim}
from scipy import signal
import matplotlib.pyplot as plt

R = 1000.0
C = 100e-6
sys = signal.TransferFunction([1], [R*C, 1])
t, y = signal.step(sys)

plt.plot(t, y)
plt.xlabel("Time (s)")
plt.ylabel("Normalized output")
plt.title("RC Low-Pass Step Response")
plt.grid(True)
plt.show()
\end{verbatim}

ผลลัพธ์ควรมีแนวโน้มเข้าใกล้ 1 เมื่อเวลามากขึ้น ซึ่งสอดคล้องกับคำตอบทางทฤษฎี

\[
y(t)=1-e^{-t/(RC)},\qquad t\ge0
\]

\subsubsection{7.4.4 SymPy: การคำนวณเชิงสัญลักษณ์}

\begin{verbatim}
import sympy as sp

t, s, R, C = sp.symbols('t s R C', positive=True)
f = sp.exp(-t/(R*C))
F = sp.laplace_transform(f, t, s, noconds=True)
print(sp.simplify(F))
\end{verbatim}

ผลเชิงสัญลักษณ์ควรสอดคล้องกับ

\[
\mathcal{L}\left\{e^{-t/(RC)}\right\}
=\frac{1}{s+1/(RC)}
\]

หรือรูปที่สมมูลกัน

\[
\frac{RC}{RCs+1}
\]

\textbf{ข้อควรระวัง:} ผลจาก symbolic software อาจอยู่ในรูปที่ต่างจากคำตอบในตำรา แต่สมมูลกันทางพีชคณิต จึงควรใช้ \texttt{simplify}, \texttt{factor} หรือการแทนค่าตรวจสอบ

\par\noindent\rule{\textwidth}{0.4pt}\par

\subsection{7.5 การใช้ MATLAB หรือ GNU Octave สำหรับงานคณิตศาสตร์วิศวกรรม}

MATLAB และ GNU Octave ใช้แนวคิดเมทริกซ์และเวกเตอร์เป็นศูนย์กลาง จึงเหมาะกับงานวิศวกรรมที่ต้องจัดการข้อมูลจำนวนมาก สมการเชิงเส้น รูปคลื่น และการวิเคราะห์ความถี่ ไวยากรณ์พื้นฐานหลายส่วนคล้ายกัน แต่ความพร้อมของ toolbox/package และรายละเอียดของฟังก์ชันอาจแตกต่างกัน

\subsubsection{7.5.1 การสร้างเวกเตอร์เวลาและรูปคลื่น}

\begin{verbatim}
fs = 5000;
T = 0.1;
t = 0:1/fs:T-1/fs;
v = 311*sin(2*pi*50*t);

plot(t,v);
xlabel('Time (s)');
ylabel('Voltage (V)');
title('50-Hz Sinusoidal Voltage');
grid on;
\end{verbatim}

\subsubsection{7.5.2 การคำนวณจำนวนเชิงซ้อน}

\begin{verbatim}
R = 20;
L = 50e-3;
C = 100e-6;
f = 50;
w = 2*pi*f;

Z = R + 1i*(w*L - 1/(w*C));
Zmag = abs(Z);
Zang = angle(Z)*180/pi;
\end{verbatim}

\subsubsection{7.5.3 การใช้ FFT}

\begin{verbatim}
Y = fft(v);
N = length(v);
faxis = (0:N-1)*(fs/N);
mag = abs(Y)/N;
\end{verbatim}

สำหรับสัญญาณจริง หากต้องการ single-sided spectrum ต้องจัดการการเลือกครึ่งสเปกตรัมและสเกลแอมพลิจูดอย่างถูกต้อง ไม่ควรนำ \texttt{abs(fft(v))} ไปอ่านเป็นขนาดแอมพลิจูดของฮาร์มอนิกทันทีโดยไม่พิจารณา normalization

\subsubsection{7.5.4 MATLAB กับ GNU Octave}

\begin{table}[htbp]
    \centering
    \begin{tabular}{l l l}
        \hline
        \textbf{ประเด็น} & \textbf{MATLAB} & \textbf{GNU Octave} \\
        \hline
        ลักษณะ & ซอฟต์แวร์เชิงพาณิชย์ & ซอฟต์แวร์เสรี/โอเพนซอร์ส \\
        รูปแบบคำสั่งพื้นฐาน & Matrix-oriented & คล้าย MATLAB ในงานพื้นฐานจำนวนมาก \\
        Toolbox/Package & มี toolbox จำนวนมาก & ใช้ core functions และ packages \\
        เหมาะกับการเรียน & ดี โดยเฉพาะเมื่อสถาบันมี license & ดีเมื่อเน้นเครื่องมือเปิดและเข้าถึงได้ \\
        ข้อควรระวัง & ฟังก์ชันบางชนิดต้องมี toolbox & package บางส่วนต้องติดตั้งเพิ่ม และความเข้ากันได้ไม่สมบูรณ์ทุกกรณี \\
        \hline
    \end{tabular}
\end{table}

\textbf{แนวทางสอนที่แนะนำ:} ใช้โจทย์เดียวกันใน Python และ MATLAB/Octave อย่างน้อยหนึ่งครั้ง เพื่อให้นักศึกษาเห็นว่าหลักคณิตศาสตร์เหมือนเดิม แม้ syntax ต่างกัน

\par\noindent\rule{\textwidth}{0.4pt}\par

\subsection{7.6 การใช้ GeoGebra หรือเครื่องมือกราฟออนไลน์ในการสร้างความเข้าใจ}

GeoGebra เหมาะกับการสำรวจแนวคิดแบบโต้ตอบ โดยเฉพาะเมื่อผู้เรียนต้องการเห็นผลของพารามิเตอร์ทันที เช่น เปลี่ยน $R$ หรือ $C$ แล้วกราฟ RC เปลี่ยนอย่างไร หรือเพิ่มจำนวนพจน์ในอนุกรมฟูเรียร์แล้วรูปคลื่นเข้าใกล้สี่เหลี่ยมอย่างไร

\subsubsection{7.6.1 ตัวอย่างการสำรวจ RC response}

กำหนด

\[
v_C(t)=V_s\left(1-e^{-t/(RC)}\right)
\]

สามารถสร้าง sliders สำหรับ $R$, $C$ และ $V_s$ แล้วสังเกตว่า time constant

\[
\tau=RC
\]

ส่งผลต่อความเร็วของการเปลี่ยนแปลงอย่างไร

\subsubsection{7.6.2 ตัวอย่างการสำรวจ Fourier partial sum}

สำหรับ square wave แบบ odd symmetry สามารถกำหนด

\[
S_N(t)=\frac{4A}{\pi}\sum_{k=0}^{N-1}\frac{1}{2k+1}\sin\left[(2k+1)\omega_0t\right]
\]

แล้วใช้ slider ของ $N$ เพื่อดูผลของจำนวนฮาร์มอนิก

% [รูปที่ 4: หน้าจอแนวคิด GeoGebra สำหรับสำรวจ RC และ Fourier]
% ประเภทภาพ: Technical Illustration / UI Concept
% ตำแหน่งที่แนะนำ: หัวข้อ 7.6
% วัตถุประสงค์การเรียนรู้: แสดงแนวคิด slider-based exploration โดยไม่ผูกกับหน้าตาเวอร์ชันใดเวอร์ชันหนึ่ง
% คำอธิบายภาพ: สองพาเนล ด้านซ้าย RC charging curve พร้อม sliders R และ C; ด้านขวา square-wave Fourier partial sum พร้อม slider N
% องค์ประกอบที่ต้องมี: กราฟ, sliders, equation field, labels tau=RC and N harmonics
% ภาษาของป้ายกำกับ: ไทยและอังกฤษ
% สัดส่วนภาพ: 16:9
% Image Generation Prompt: Two-panel educational math software concept. Left panel: interactive RC charging curve vC(t)=Vs(1-exp(-t/RC)) with sliders for R and C and displayed time constant tau=RC. Right panel: Fourier partial-sum approximation of a square wave with slider N controlling the number of odd harmonics. Clean white background, bilingual Thai-English labels, precise engineering textbook vector style, 16:9.
% Alt Text: ภาพแนวคิดการใช้ slider สำรวจผลของ R, C และจำนวนฮาร์มอนิกต่อกราฟ

\begin{figure}[htbp]
    \centering
    \includegraphics[width=0.8\textwidth]{images/figure-04.png}
    \caption{หน้าจอแนวคิด GeoGebra สำหรับสำรวจ RC และ Fourier}
\end{figure}

\par\noindent\rule{\textwidth}{0.4pt}\par
\subsection{7.7 การคำนวณจำนวนเชิงซ้อนและเฟสเซอร์ด้วยโปรแกรม}

\subsubsection{7.7.1 จากจำนวนเชิงซ้อนสู่เฟสเซอร์}

จำนวนเชิงซ้อนในรูปสี่เหลี่ยมเขียนได้เป็น

\[
Z=a+jb
\]

และแปลงเป็นรูปเชิงขั้วได้เป็น

\[
Z=|Z|\angle\theta
\]

เมื่อ

\[
|Z|=\sqrt{a^2+b^2}
\]

และ

\[
\theta=\operatorname{atan2}(b,a)
\]

การใช้ \texttt{atan2} มีความสำคัญกว่าการใช้เพียง $\tan^{-1}(b/a)$ เพราะสามารถระบุควอดแรนต์ของมุมได้ถูกต้อง

\subsubsection{7.7.2 สมการอิมพีแดนซ์พื้นฐาน}

สำหรับวงจรไซน์คงตัวที่ความถี่เชิงมุม $\omega$

\[
Z_R=R
\]

\[
Z_L=j\omega L
\]

\[
Z_C=\frac{1}{j\omega C}=-j\frac{1}{\omega C}
\]

โดย

\begin{itemize}
    \item $R$ หน่วยโอห์ม $(\Omega)$
    \item $L$ หน่วยเฮนรี (H)
    \item $C$ หน่วยฟารัด (F)
    \item $\omega=2\pi f$ หน่วย rad/s
\end{itemize}

สมการเหล่านี้ใช้กับการวิเคราะห์ sinusoidal steady state และสมมติว่าระบบเชิงเส้นที่ความถี่เดียวกัน

\subsubsection{7.7.3 การคำนวณอิมพีแดนซ์เป็นฟังก์ชันของความถี่}

พิจารณาวงจร RLC อนุกรม

\[
Z(f)=R+j\left(2\pi fL-\frac{1}{2\pi fC}\right)
\]

เรโซแนนซ์เกิดเมื่อส่วนจินตภาพเป็นศูนย์

\[
2\pi f_0L=\frac{1}{2\pi f_0C}
\]

จึงได้

\[
f_0=\frac{1}{2\pi\sqrt{LC}}
\]

สำหรับ $L=50\ \text{mH}$ และ $C=100\ \mu\text{F}$

\[
f_0\approx\frac{1}{2\pi\sqrt{(0.05)(100\times10^{-6})}}
\approx71.18\ \text{Hz}
\]

ที่ความถี่นี้อิมพีแดนซ์ในอุดมคติของวงจรอนุกรมจะมีส่วนรีแอกทีฟใกล้ศูนย์และ $|Z|$ ต่ำสุดใกล้ $R$

\begin{verbatim}
import numpy as np
import matplotlib.pyplot as plt

R = 20.0
L = 50e-3
C = 100e-6

f = np.logspace(0, 3, 800)
w = 2*np.pi*f
Z = R + 1j*(w*L - 1/(w*C))

plt.semilogx(f, np.abs(Z))
plt.xlabel("Frequency (Hz)")
plt.ylabel("|Z| (ohm)")
plt.title("Series RLC Impedance Magnitude")
plt.grid(True, which="both")
plt.show()
\end{verbatim}

% [รูปที่ 5: ขนาดอิมพีแดนซ์และมุมเฟสของวงจร RLC ตามความถี่]
% ประเภทภาพ: Graph
% ตำแหน่งที่แนะนำ: หัวข้อ 7.7.3
% วัตถุประสงค์การเรียนรู้: แสดงการเชื่อมโยงจำนวนเชิงซ้อนกับพฤติกรรมความถี่ของวงจร
% คำอธิบายภาพ: กราฟสองส่วนแต่ไม่จำเป็นต้องอยู่ใน subplot เดียว: magnitude $|Z|$ และ phase $\angle Z$ เทียบความถี่แบบ logarithmic โดยทำเครื่องหมาย $f_0\approx71.18$ Hz
% องค์ประกอบที่ต้องมี: แกนความถี่ Hz, magnitude ohm, phase degree, resonance marker
% ภาษาของป้ายกำกับ: ไทยและอังกฤษ
% สัดส่วนภาพ: 16:9
% Graph Generation Prompt: Plot the series RLC impedance Z(f)=20+j(2πf·0.05−1/(2πf·100e−6)) for f from 1 Hz to 1000 Hz on logarithmic frequency axes. Show magnitude |Z| in ohms and phase angle in degrees, clearly mark resonance f0≈71.18 Hz where the imaginary part crosses zero. Use precise axes, units, legends, white engineering textbook background, bilingual Thai-English labels, 16:9.
% Alt Text: กราฟอิมพีแดนซ์ของวงจร RLC แสดงค่าต่ำสุดและมุมใกล้ศูนย์ที่ความถี่เรโซแนนซ์

\begin{figure}[htbp]
    \centering
    \includegraphics[width=0.8\textwidth]{images/figure-05.png}
    \caption{ขนาดอิมพีแดนซ์และมุมเฟสของวงจร RLC ตามความถี่}
\end{figure}

\subsubsection{7.7.4 การตรวจสอบเฟสเซอร์ด้วยหลักทางกายภาพ}

ก่อนยอมรับผลลัพธ์จากโปรแกรม ควรตั้งคำถามต่อไปนี้

\begin{itemize}
    \item วงจรมีลักษณะ inductive หรือ capacitive?
    \item กระแสควรนำหรือล้าหลังแรงดัน?
    \item ขนาดกระแสควรอยู่ในช่วงใดจาก $I\approx V/|Z|$?
    \item หาก $f$ เพิ่มมากขึ้น $X_L$ และ $X_C$ เปลี่ยนอย่างไร?
\end{itemize}

คำถามเหล่านี้เป็น \textbf{engineering sanity check} ที่ช่วยจับข้อผิดพลาดของหน่วยและเครื่องหมาย

\par\noindent\rule{\textwidth}{0.4pt}\par

\subsection{7.8 การวิเคราะห์เวกเตอร์ด้วยโปรแกรม}

\subsubsection{7.8.1 การคำนวณเวกเตอร์เชิงตัวเลข}

ให้เวกเตอร์

\[
\mathbf{A}=2\hat{\mathbf{x}}-\hat{\mathbf{y}}+3\hat{\mathbf{z}}
\]

และ

\[
\mathbf{B}=4\hat{\mathbf{x}}+2\hat{\mathbf{y}}-\hat{\mathbf{z}}
\]

ดอตโปรดักต์คือ

\[
\mathbf{A}\cdot\mathbf{B}=2(4)+(-1)(2)+3(-1)=3
\]

ครอสโปรดักต์คือ

\[
\mathbf{A}\times\mathbf{B}
=\begin{vmatrix}
\hat{\mathbf{x}} & \hat{\mathbf{y}} & \hat{\mathbf{z}}\\
2&-1&3\\
4&2&-1
\end{vmatrix}
\]

คำนวณด้วย NumPy ได้ดังนี้

\begin{verbatim}
import numpy as np

A = np.array([2.0, -1.0, 3.0])
B = np.array([4.0,  2.0, -1.0])

dot_AB = np.dot(A, B)
cross_AB = np.cross(A, B)

print(dot_AB)
print(cross_AB)
\end{verbatim}

การตรวจสอบที่ทำได้ทันทีคือ

\[
(\mathbf{A}\times\mathbf{B})\cdot\mathbf{A}=0
\]

และ

\[
(\mathbf{A}\times\mathbf{B})\cdot\mathbf{B}=0
\]

เพราะครอสโปรดักต์ต้องตั้งฉากกับทั้งสองเวกเตอร์

\subsubsection{7.8.2 เกรเดียนต์เชิงสัญลักษณ์}

ให้ศักย์ไฟฟ้าเชิงตัวอย่าง

\[
V(x,y,z)=3x^2y+2yz^2\quad \text{V}
\]

เกรเดียนต์คือ

\[
\nabla V=
\frac{\partial V}{\partial x}\hat{\mathbf{x}}+
\frac{\partial V}{\partial y}\hat{\mathbf{y}}+
\frac{\partial V}{\partial z}\hat{\mathbf{z}}
\]

ดังนั้น

\[
\nabla V=
6xy\hat{\mathbf{x}}+
(3x^2+2z^2)\hat{\mathbf{y}}+
4yz\hat{\mathbf{z}}
\]

สนามไฟฟ้าสัมพันธ์กับศักย์โดย

\[
\mathbf{E}=-\nabla V
\]

\textbf{ความหมายทางกายภาพ:} เครื่องหมายลบหมายถึงสนามไฟฟ้าชี้ไปในทิศที่ศักย์ลดลงเร็วที่สุด

\begin{verbatim}
import sympy as sp

x, y, z = sp.symbols('x y z', real=True)
V = 3*x**2*y + 2*y*z**2

gradV = sp.Matrix([
    sp.diff(V, x),
    sp.diff(V, y),
    sp.diff(V, z)
])

E = -gradV
print(E)
\end{verbatim}

\subsubsection{7.8.3 ไดเวอร์เจนซ์และเคิร์ล}

สำหรับสนามเวกเตอร์

\[
\mathbf{F}=F_x\hat{\mathbf{x}}+F_y\hat{\mathbf{y}}+F_z\hat{\mathbf{z}}
\]

ไดเวอร์เจนซ์คือ

\[
\nabla\cdot\mathbf{F}
=\frac{\partial F_x}{\partial x}
+\frac{\partial F_y}{\partial y}
+\frac{\partial F_z}{\partial z}
\]

และเคิร์ลคือ

\[
\nabla\times\mathbf{F}
=
\begin{vmatrix}
\hat{\mathbf{x}}&\hat{\mathbf{y}}&\hat{\mathbf{z}}\\
\partial/\partial x&\partial/\partial y&\partial/\partial z\\
F_x&F_y&F_z
\end{vmatrix}
\]

การใช้ symbolic software เหมาะเมื่อส่วนประกอบของสนามเป็นฟังก์ชันของหลายตัวแปร แต่ผู้เรียนควรตรวจหนึ่งองค์ประกอบด้วยมือก่อนเพื่อป้องกันการกำหนดตัวแปรผิด

% [รูปที่ 6: การวิเคราะห์สนามเวกเตอร์จากศักย์ด้วยโปรแกรม]
% ประเภทภาพ: Concept Diagram / Vector Field
% ตำแหน่งที่แนะนำ: หัวข้อ 7.8
% วัตถุประสงค์การเรียนรู้: เชื่อมโยงสมการศักย์ เกรเดียนต์ และทิศสนามไฟฟ้า
% คำอธิบายภาพ: แสดงพื้นผิวศักย์แบบ contour/3D ทางซ้าย และลูกศรสนาม $\mathbf{E}=-\nabla V$ ทางขวา โดยลูกศรชี้จากศักย์สูงไปต่ำ
% องค์ประกอบที่ต้องมี: equipotential contours, gradient arrow, electric-field arrow, coordinate axes
% ภาษาของป้ายกำกับ: ไทยและอังกฤษ
% สัดส่วนภาพ: 16:9
% Image Generation Prompt: Engineering mathematics illustration connecting electric potential and vector field. Show equipotential contours or a smooth potential surface V(x,y) and arrows for grad V pointing toward steepest increase, with electric field E=-grad V pointing in the opposite direction toward lower potential. Include coordinate axes and concise bilingual Thai-English labels. White background, clean vector textbook style, 16:9.
% Alt Text: ภาพแสดงเกรเดียนต์ของศักย์และสนามไฟฟ้าที่ชี้ทิศตรงข้ามกับเกรเดียนต์

\begin{figure}[htbp]
    \centering
    \includegraphics[width=0.8\textwidth]{images/figure-06.png}
    \caption{การวิเคราะห์สนามเวกเตอร์จากศักย์ด้วยโปรแกรม}
\end{figure}

\par\noindent\rule{\textwidth}{0.4pt}\par

\subsection{7.9 การสร้างและวิเคราะห์อนุกรมฟูเรียร์ด้วยโปรแกรม}

\subsubsection{7.9.1 จุดเชื่อมระหว่าง Fourier series และ FFT}

อนุกรมฟูเรียร์เป็นแบบจำลองเชิงคณิตศาสตร์ของฟังก์ชันเป็นคาบ ส่วน FFT เป็นอัลกอริทึมสำหรับคำนวณ DFT ของข้อมูลที่สุ่มตัวอย่างแล้ว ทั้งสองแนวคิดเกี่ยวข้องกับองค์ประกอบความถี่ แต่ไม่ควรถือว่าเหมือนกันทุกประการ

สำหรับข้อมูล $x[n]$ จำนวน $N$ จุด DFT นิยามได้เป็น

\[
X[k]=\sum_{n=0}^{N-1}x[n]e^{-j2\pi kn/N}
\]

เมื่อ

\begin{itemize}
    \item $n$ คือดัชนีตัวอย่างในโดเมนเวลา
    \item $k$ คือดัชนี bin ในโดเมนความถี่
    \item $N$ คือจำนวนตัวอย่าง
    \item $X[k]$ เป็นจำนวนเชิงซ้อน แสดงทั้งขนาดและเฟส
\end{itemize}

ถ้า sampling frequency เท่ากับ $f_s$ แกนความถี่สัมพันธ์กับ bin โดย

\[
f_k=\frac{k f_s}{N}
\]

และความละเอียดเชิงความถี่คือ

\[
\Delta f=\frac{f_s}{N}=\frac{1}{T_{\text{record}}}
\]

\subsubsection{7.9.2 เงื่อนไข Nyquist แบบพื้นฐาน}

เพื่อหลีกเลี่ยง aliasing สำหรับสัญญาณที่มีองค์ประกอบความถี่สูงสุด $f_{\max}$ ควรมี

\[
f_s>2f_{\max}
\]

ในทางปฏิบัติงานจริงมักต้องใช้ sampling rate สูงกว่าค่าขั้นต่ำและมี anti-aliasing filter ก่อน ADC แต่ในบทนี้เน้นแนวคิดทางคณิตศาสตร์

\subsubsection{ตัวอย่างที่ 3 FFT ของสัญญาณที่มี 50 Hz และ 150 Hz}

\textbf{ตัวอย่างเพิ่มเติมเพื่อการเรียนรู้}

กำหนด

\[
x(t)=10\sin(2\pi50t)+3\sin(2\pi150t)
\]

สุ่มด้วย $f_s=2000$ Hz เป็นเวลา 1 s

\begin{verbatim}
import numpy as np
import matplotlib.pyplot as plt

fs = 2000.0
T = 1.0
t = np.arange(0, T, 1/fs)

x = 10*np.sin(2*np.pi*50*t) + 3*np.sin(2*np.pi*150*t)

N = len(x)
X = np.fft.rfft(x)
freq = np.fft.rfftfreq(N, d=1/fs)

amp = np.abs(X) / N
amp[1:-1] *= 2

plt.stem(freq, amp)
plt.xlim(0, 300)
plt.xlabel("Frequency (Hz)")
plt.ylabel("Amplitude")
plt.title("Single-Sided Amplitude Spectrum")
plt.grid(True)
plt.show()
\end{verbatim}

\textbf{ผลที่คาดหวัง:} เห็นยอดเด่นใกล้ 50 Hz ที่แอมพลิจูดประมาณ 10 และ 150 Hz ที่แอมพลิจูดประมาณ 3

\textbf{เหตุผลที่เลือก $T=1$ s:}

\[
\Delta f=\frac{1}{T}=1\ \text{Hz}
\]

ดังนั้น 50 และ 150 Hz ตรงกับ FFT bins พอดีในกรณีอุดมคตินี้ ทำให้การอ่านแอมพลิจูดง่ายขึ้น

\subsubsection{7.9.3 spectral leakage}

ถ้าช่วงเวลาที่บันทึกไม่ครอบคลุมจำนวนคาบที่ลงตัว สัญญาณที่ปลายซ้ายและขวาของ record อาจไม่ต่อเนื่องเมื่อนำไปมองเป็นคาบ ส่งผลให้พลังงานกระจายไปหลาย bins เรียกว่า \textbf{spectral leakage}

จึงไม่ควรสรุปว่า “ยอด FFT ที่กระจายแปลว่าระบบมีหลายความถี่จริงเสมอ” ต้องพิจารณา

\begin{itemize}
    \item ความยาว record
    \item sampling rate
    \item window function
    \item noise
    \item frequency resolution
\end{itemize}

\subsubsection{7.9.4 การสร้าง Fourier partial sum ของ square wave}

สำหรับ square wave แอมพลิจูด $\pm A$ ที่มี odd symmetry

\[
x(t)=\frac{4A}{\pi}\left[
\sin(\omega_0t)+\frac{1}{3}\sin(3\omega_0t)+\frac{1}{5}\sin(5\omega_0t)+\cdots
\right]
\]

ตัวอย่าง Python

\begin{verbatim}
import numpy as np
import matplotlib.pyplot as plt

A = 1.0
f0 = 50.0
w0 = 2*np.pi*f0

t = np.linspace(0, 2/f0, 2000, endpoint=False)

for N in [1, 3, 10, 50]:
    y = np.zeros_like(t)
    for k in range(N):
        n = 2*k + 1
        y += (4*A/np.pi) * np.sin(n*w0*t) / n
    plt.plot(t, y, label=f"N={N}")

plt.xlabel("Time (s)")
plt.ylabel("Amplitude")
plt.title("Square-Wave Fourier Partial Sums")
plt.legend()
plt.grid(True)
plt.show()
\end{verbatim}

ผู้เรียนควรสังเกตว่าเมื่อ $N$ เพิ่มขึ้น รูปคลื่นเข้าใกล้ square wave มากขึ้น แต่บริเวณ discontinuity ยังคงเกิด overshoot ในลักษณะ Gibbs phenomenon

% [รูปที่ 7: จากสัญญาณโดเมนเวลาสู่ FFT spectrum]
% ประเภทภาพ: Three-tier Engineering Illustration
% ตำแหน่งที่แนะนำ: หัวข้อ 7.9
% วัตถุประสงค์การเรียนรู้: เชื่อมเวลา การสุ่มตัวอย่าง FFT และสเปกตรัม
% คำอธิบายภาพ: ชั้นบนสัญญาณ $x(t)$ ที่มี 50 Hz และ 150 Hz; ชั้นกลาง sampled points พร้อม $f_s$; ชั้นล่าง single-sided spectrum ที่มียอด 50 และ 150 Hz
% องค์ประกอบที่ต้องมี: time axis, sample markers, frequency axis, FFT arrows, peaks
% ภาษาของป้ายกำกับ: ไทยและอังกฤษ
% สัดส่วนภาพ: 16:9
% Image Generation Prompt: Three-tier engineering signal illustration. Top: continuous waveform x(t)=10 sin(2π50t)+3 sin(2π150t). Middle: uniformly sampled points at fs=2000 Hz with sample interval indicated. Bottom: single-sided FFT amplitude spectrum with clear peaks at 50 Hz amplitude 10 and 150 Hz amplitude 3. Show data-flow arrows from time signal to samples to FFT spectrum. Clean white-background textbook vector style, bilingual Thai-English labels, 16:9.
% Graph Generation Prompt: Generate x(t)=10 sin(2π·50 t)+3 sin(2π·150 t) for 0≤t<1 s, fs=2000 Hz. Plot a short time segment and a single-sided FFT amplitude spectrum from 0–300 Hz. Mark 50 Hz and 150 Hz peaks and label units.
% Alt Text: ภาพสามระดับแสดงสัญญาณเวลา การสุ่มตัวอย่าง และสเปกตรัม FFT ที่มีองค์ประกอบ 50 และ 150 เฮิรตซ์

\begin{figure}[htbp]
    \centering
    \includegraphics[width=0.8\textwidth]{images/figure-07.png}
    \caption{จากสัญญาณโดเมนเวลาสู่ FFT spectrum}
\end{figure}

\par\noindent\rule{\textwidth}{0.4pt}\par

\subsection{7.10 การคำนวณการแปลงลาปลาซและการวิเคราะห์วงจรด้วยโปรแกรม}

\subsubsection{7.10.1 นิยามและขอบเขต}

การแปลงลาปลาซแบบ unilateral ที่ใช้ในงานวงจรเขียนเป็น

\[
F(s)=\mathcal{L}\{f(t)\}
=\int_{0^-}^{\infty} f(t)e^{-st}\,dt
\]

เมื่อ

\[
s=\sigma+j\omega
\]

การใช้โปรแกรม symbolic เหมาะกับ

\begin{itemize}
    \item หา Laplace transform และ inverse Laplace transform
    \item จัดรูปเศษส่วนย่อย
    \item ตรวจสมการคำตอบจากการแก้มือ
    \item แก้สมการเชิงอนุพันธ์พร้อมเงื่อนไขเริ่มต้นในกรณีที่รองรับ
\end{itemize}

\subsubsection{7.10.2 ตัวอย่างพื้นฐาน: Laplace ของสัญญาณเอ็กซ์โพเนนเชียล}

ให้

\[
f(t)=e^{-at}u(t),\qquad a>0
\]

จากนิยาม

\[
F(s)=\int_0^\infty e^{-at}e^{-st}\,dt
=\int_0^\infty e^{-(s+a)t}\,dt
\]

จึงได้

\[
F(s)=\frac{1}{s+a}
\]

โดยมีเงื่อนไขการลู่เข้าที่เกี่ยวข้องกับส่วนจริงของ $s$

SymPy

\begin{verbatim}
import sympy as sp

t, s, a = sp.symbols('t s a', positive=True)
f = sp.exp(-a*t)
F = sp.laplace_transform(f, t, s, noconds=True)
print(F)
\end{verbatim}

MATLAB Symbolic

\begin{verbatim}
syms t s a positive
f = exp(-a*t);
F = laplace(f,t,s);
\end{verbatim}

\subsubsection{7.10.3 การวิเคราะห์ RC low-pass}

พิจารณาวงจร RC ที่รับแรงดันอินพุต $v_i(t)$ และวัดเอาต์พุตที่ตัวเก็บประจุ $v_C(t)$

จาก KVL

\[
v_i(t)=Ri(t)+v_C(t)
\]

และ

\[
i(t)=C\frac{dv_C(t)}{dt}
\]

จึงได้

\[
RC\frac{dv_C(t)}{dt}+v_C(t)=v_i(t)
\]

ถ้าเงื่อนไขเริ่มต้นเป็นศูนย์และแปลงลาปลาซ

\[
(RCs+1)V_C(s)=V_i(s)
\]

ดังนั้นฟังก์ชันถ่ายโอนคือ

\[
H(s)=\frac{V_C(s)}{V_i(s)}=\frac{1}{RCs+1}
\]

\textbf{สมมติฐาน:} อุปกรณ์เป็นอุดมคติ ระบบเป็นเชิงเส้น และนิยาม transfer function ใช้ zero initial conditions

\subsubsection{ตัวอย่างที่ 4 เปรียบเทียบ symbolic solution กับ numerical step response}

\textbf{ตัวอย่างเพิ่มเติมเพื่อการเรียนรู้}

กำหนด $R=1\ \text{k}\Omega$, $C=100\ \mu\text{F}$ และ step input ขนาด 5 V

time constant คือ

\[
\tau=RC=(1000)(100\times10^{-6})=0.1\ \text{s}
\]

คำตอบเชิงทฤษฎี

\[
v_C(t)=5\left(1-e^{-t/0.1}\right)u(t)
\]

ที่ $t=\tau=0.1$ s

\[
v_C(\tau)=5(1-e^{-1})\approx3.161\ \text{V}
\]

หรือประมาณ 63.2\% ของค่าสุดท้าย

\textbf{Python: symbolic + numerical comparison}

\begin{verbatim}
import numpy as np
import sympy as sp
import matplotlib.pyplot as plt
from scipy import signal

# symbolic
s, t = sp.symbols('s t', positive=True)
R = 1000.0
C = 100e-6
Vs = 5.0

H = 1/(R*C*s + 1)
Vc_s = H * Vs/s
Vc_t = sp.inverse_laplace_transform(Vc_s, s, t)
print(sp.simplify(Vc_t))

# numerical LTI step response (unit step, then scale by Vs)
sys = signal.TransferFunction([1], [R*C, 1])
t_num = np.linspace(0, 0.5, 500)
t_out, y_unit = signal.step(sys, T=t_num)
y_num = Vs*y_unit

# analytical numeric curve
v_exact = Vs*(1 - np.exp(-t_num/(R*C)))

plt.plot(t_out, y_num, label="SciPy step")
plt.plot(t_num, v_exact, "--", label="Analytical")
plt.xlabel("Time (s)")
plt.ylabel("Capacitor voltage (V)")
plt.title("RC Step Response: Numerical vs Analytical")
plt.legend()
plt.grid(True)
plt.show()
\end{verbatim}

\textbf{การตรวจสอบ 4 จุด}

\begin{enumerate}
    \item $v_C(0^+)\approx0$ V เมื่อเริ่มจากตัวเก็บประจุไม่มีประจุ
    \item $v_C(\tau)\approx3.16$ V
    \item $v_C(5\tau)$ ควรใกล้ 5 V มาก
    \item กราฟ numerical และ analytical ควรทับกันเกือบทั้งหมดในกรณีนี้
\end{enumerate}

\subsubsection{7.10.4 ฟังก์ชันถ่ายโอน โพล และความหมายเชิงกายภาพ}

สำหรับ

\[
H(s)=\frac{1}{RCs+1}
\]

โพลคือรากของ denominator

\[
RCs+1=0
\]

จึงได้

\[
p=-\frac{1}{RC}
\]

เมื่อ $R>0$ และ $C>0$ โพลอยู่ทางซ้ายของระนาบ $s$ สอดคล้องกับการตอบสนองที่ลดทอนและระบบอันดับหนึ่งที่เสถียรภายใต้แบบจำลองนี้

โปรแกรมสามารถหาค่าโพลได้ง่าย แต่ผู้เรียนต้องตีความว่า \textbf{ตำแหน่งโพลสัมพันธ์กับอัตราการเปลี่ยนแปลงของ transient} ไม่ใช่เพียงตัวเลขที่โปรแกรมแสดง

% [รูปที่ 8: กระบวนการวิเคราะห์วงจร RC จากโดเมนเวลาไปยังผลตอบสนอง]
% ประเภทภาพ: Flowchart / Block Diagram
% ตำแหน่งที่แนะนำ: หัวข้อ 7.10.3
% วัตถุประสงค์การเรียนรู้: เชื่อม KVL → ODE → Laplace → transfer function → pole → step response
% คำอธิบายภาพ: วงจร RC low-pass ทางซ้าย ต่อด้วยสมการ $RC\dot v_C+v_C=v_i$, แปลงเป็น $H(s)=1/(RCs+1)$, แสดง pole $-1/RC$, และกราฟ step response
% องค์ประกอบที่ต้องมี: R, C, input/output, differential equation, s-domain, pole marker, exponential response
% ภาษาของป้ายกำกับ: ไทยและอังกฤษ
% สัดส่วนภาพ: 16:9
% Image Generation Prompt: Engineering education flow diagram for an RC low-pass circuit. Show the physical RC circuit with input vi(t) and capacitor output vC(t), derive RC dvC/dt + vC = vi, transform to H(s)=1/(RCs+1), show pole p=-1/RC on the negative real s-axis, and show the exponential step response vC(t)=Vs(1-exp(-t/RC)). Bilingual Thai-English labels, clean white textbook vector style, 16:9.
% Alt Text: แผนภาพเชื่อมวงจร RC สมการเชิงอนุพันธ์ ฟังก์ชันถ่ายโอน โพล และกราฟการตอบสนองขั้นบันได

\begin{figure}[htbp]
    \centering
    \includegraphics[width=0.8\textwidth]{images/figure-08.png}
    \caption{กระบวนการวิเคราะห์วงจร RC จากโดเมนเวลาไปยังผลตอบสนอง}
\end{figure}

\subsubsection{7.10.5 ตัวอย่างประยุกต์: วงจร RLC อันดับสอง}

วงจร RLC อนุกรมที่วัดแรงดันตัวเก็บประจุสามารถมีฟังก์ชันถ่ายโอนรูป

\[
H(s)=\frac{1/LC}{s^2+(R/L)s+1/LC}
\]

เทียบกับรูปมาตรฐานอันดับสอง

\[
H(s)=\frac{\omega_n^2}{s^2+2\zeta\omega_ns+\omega_n^2}
\]

จึงระบุได้ว่า

\[
\omega_n=\frac{1}{\sqrt{LC}}
\]

และ

\[
2\zeta\omega_n=\frac{R}{L}
\]

หรือ

\[
\zeta=\frac{R}{2}\sqrt{\frac{C}{L}}
\]

โปรแกรมช่วย sweep ค่า $R$ เพื่อดูการเปลี่ยนจาก underdamped ไป critical/overdamped ได้อย่างชัดเจน

\begin{verbatim}
import numpy as np
import matplotlib.pyplot as plt
from scipy import signal

L = 50e-3
C = 100e-6
R_values = [5, 20, 50]

t = np.linspace(0, 0.2, 1200)

for R in R_values:
    num = [1/(L*C)]
    den = [1, R/L, 1/(L*C)]
    sys = signal.TransferFunction(num, den)
    tout, y = signal.step(sys, T=t)
    plt.plot(tout, y, label=f"R={R} ohm")

plt.xlabel("Time (s)")
plt.ylabel("Normalized capacitor voltage")
plt.title("RLC Step Response for Different R")
plt.legend()
plt.grid(True)
plt.show()
\end{verbatim}

\textbf{สิ่งที่ควรสังเกต:} เมื่อ $R$ เพิ่ม ค่า damping ratio เพิ่มขึ้น การสั่นลดลง แต่ระบบอาจตอบสนองช้าลง จึงต้องตีความจากทั้ง overshoot และ settling behavior

\par\noindent\rule{\textwidth}{0.4pt}\par

\subsection{7.11 การจัดทำรายงานผลการคำนวณทางวิศวกรรม}

ผลลัพธ์ที่ดีไม่ใช่เพียงรูปกราฟสวยหรือคำตอบสุดท้าย แต่ต้องทำให้ผู้อ่านตรวจสอบที่มาได้ รายงานการคำนวณอย่างน้อยควรมีองค์ประกอบต่อไปนี้

\subsubsection{7.11.1 โครงสร้างรายงานที่แนะนำ}

\begin{enumerate}
    \item \textbf{ชื่อปัญหาและวัตถุประสงค์}
    ระบุว่าต้องการคำนวณหรือศึกษาพฤติกรรมใด
    \item \textbf{แบบจำลองและสมมติฐาน}
    แสดงวงจร สมการ และข้อสมมติ เช่น linear, ideal component, zero initial condition
    \item \textbf{ข้อมูลนำเข้าและหน่วย}
    จัดตาราง $R$, $L$, $C$, $V$, $f$, $f_s$, $T$ เป็นต้น
    \item \textbf{วิธีคำนวณ}
    ระบุเครื่องมือ ไลบรารี ฟังก์ชัน และขั้นตอนที่สำคัญ
    \item \textbf{ผลลัพธ์}
    แสดงค่าตัวเลข ตาราง และกราฟพร้อมหน่วย
    \item \textbf{การตรวจสอบ}
    เปรียบเทียบกับค่าจากมือ/สมการวิเคราะห์/พฤติกรรมปลายขอบ
    \item \textbf{อภิปรายผล}
    อธิบายสิ่งที่กราฟหรือค่าคำนวณบอกเกี่ยวกับระบบ
    \item \textbf{ข้อจำกัด}
    เช่น sampling, ideal model, numerical precision, missing parasitics
    \item \textbf{สรุป}
    ตอบวัตถุประสงค์โดยตรง
    \item \textbf{ไฟล์ประกอบ}
    แนบ script/notebook/data file เพื่อให้ทำซ้ำได้
\end{enumerate}

\subsubsection{7.11.2 ตารางพารามิเตอร์ตัวอย่าง}

\begin{table}[htbp]
    \centering
    \begin{tabular}{l r r l l}
        \hline
        \textbf{พารามิเตอร์} & \textbf{สัญลักษณ์} & \textbf{ค่า} & \textbf{หน่วย} & \textbf{แหล่งที่มา/หมายเหตุ} \\
        \hline
        ความต้านทาน & $R$ & 1000 & $\Omega$ & ค่ากำหนด \\
        ความจุ & $C$ & 100 & $\mu$F & ค่ากำหนด \\
        แรงดัน step & $V_s$ & 5 & V & ค่ากำหนด \\
        time constant & $\tau$ & 0.1 & s & คำนวณจาก $RC$ \\
        \hline
    \end{tabular}
\end{table}

\subsubsection{7.11.3 การรายงานกราฟ}

กราฟควรตอบคำถามได้ว่า

\begin{itemize}
    \item แกน $x$ คืออะไรและหน่วยอะไร?
    \item แกน $y$ คืออะไรและหน่วยอะไร?
    \item เส้นแต่ละเส้นแทนกรณีใด?
    \item มีพารามิเตอร์ใดคงที่?
    \item มีค่าจุดสำคัญใด เช่น $f_0$, $\tau$, peak หรือ cutoff?
    \item ข้อสรุปจากกราฟคืออะไร?
\end{itemize}

ไม่ควรใส่กราฟโดยไม่มีคำอธิบาย เช่น “จากรูปที่ 8 พบว่าเมื่อ $R$ เพิ่ม การสั่นลดลง สอดคล้องกับค่า damping ratio ที่เพิ่มขึ้น”

\subsubsection{7.11.4 การรายงานตัวเลขและ significant figures}

หากค่าพารามิเตอร์มีความแม่นยำจำกัด ไม่ควรรายงานคำตอบทศนิยมจำนวนมากเกินความจำเป็น ตัวอย่างเช่น $R=1.0\ \text{k}\Omega$ และ $C=100\ \mu\text{F}$ ไม่ได้ทำให้ $\tau=0.1000000000$ s มีความหมายทางกายภาพมากขึ้น

\subsubsection{7.11.5 หลักความสามารถทำซ้ำ}

รายงานควรทำให้ผู้ตรวจสามารถ

\begin{enumerate}
    \item ทราบสมการที่ใช้
    \item ทราบค่าพารามิเตอร์
    \item ทราบเครื่องมือหรือไลบรารี
    \item รันไฟล์ต้นฉบับได้
    \item ได้กราฟหรือคำตอบสอดคล้องกับรายงาน
\end{enumerate}

\par\noindent\rule{\textwidth}{0.4pt}\par

\subsection{7.12 กิจกรรมปฏิบัติและแบบฝึกหัดบูรณาการ}

\subsubsection{กิจกรรมที่ 1: มือ–โปรแกรม–เหตุผลทางกายภาพ}

\textbf{ผลลัพธ์การเรียนรู้ที่เกี่ยวข้อง:} CLO2, CLO5 \\
\textbf{ระยะเวลา:} 30–40 นาที \\
\textbf{รูปแบบ:} คู่ \\
\textbf{อุปกรณ์:} เครื่องคิดเลขและ Python หรือ MATLAB/Octave

\textbf{โจทย์:} ใช้วงจร RLC ในตัวอย่างที่ 1

\textbf{ขั้นตอน}

\begin{enumerate}
    \item คนที่ 1 คำนวณ $Z$ และ $I$ ด้วยเครื่องคิดเลข
    \item คนที่ 2 คำนวณด้วยโปรแกรม
    \item สลับกันตรวจหน่วยและมุมเฟส
    \item อธิบายว่ากระแสนำหรือล้าหลังแรงดันจากค่า $X_L-X_C$
    \item เปลี่ยนความถี่เป็น 100 Hz แล้วทำนาย “ก่อนคำนวณ” วงจรจะมีลักษณะ inductive หรือ capacitive
    \item ใช้โปรแกรมยืนยันการทำนาย
\end{enumerate}

\textbf{ผลลัพธ์ที่คาดหวัง:} ผู้เรียนเชื่อมคำตอบเชิงตัวเลขกับพฤติกรรมวงจร ไม่มองจำนวนเชิงซ้อนเป็นเพียงรูปแบบข้อมูล

\textbf{การประเมิน:} ความถูกต้อง 40\%, เหตุผลทางกายภาพ 40\%, การตรวจหน่วยและการสื่อสาร 20\%

\subsubsection{กิจกรรมที่ 2: Fourier Explorer}

\textbf{ผลลัพธ์การเรียนรู้ที่เกี่ยวข้อง:} CLO3, CLO5 \\
\textbf{ระยะเวลา:} 45–60 นาที \\
\textbf{รูปแบบ:} กลุ่ม 3 คน \\
\textbf{อุปกรณ์:} Python/MATLAB/Octave หรือ GeoGebra

\textbf{โจทย์:} ศึกษา square-wave Fourier partial sum ที่ $N=1,3,5,15,50$

\textbf{ภารกิจ}

\begin{itemize}
    \item สร้างกราฟทุกกรณี
    \item ระบุว่าฮาร์มอนิกใดถูกใช้
    \item อภิปรายการเปลี่ยนแปลงของรูปคลื่นเมื่อ $N$ เพิ่ม
    \item ขยายบริเวณ discontinuity และอธิบาย Gibbs phenomenon เชิงคุณภาพ
    \item เปรียบเทียบข้อดีของ GeoGebra slider กับ script
\end{itemize}

\subsubsection{กิจกรรมที่ 3: RC Computational Report}

\textbf{ผลลัพธ์การเรียนรู้ที่เกี่ยวข้อง:} CLO4, CLO5, CLO6 \\
\textbf{ระยะเวลา:} 60–90 นาทีในชั้นเรียนและเก็บงานต่อได้ตามความเหมาะสม \\
\textbf{รูปแบบ:} รายบุคคลหรือคู่

\textbf{โจทย์:} สำหรับ RC low-pass $R=2.2\ \text{k}\Omega$, $C=47\ \mu\text{F}$ และ step input 5 V

ผู้เรียนต้อง

\begin{enumerate}
    \item คำนวณ $\tau$ ด้วยมือ
    \item สร้าง $H(s)$
    \item หา step response ด้วยโปรแกรม
    \item ระบุ $v_C(\tau)$ และ $v_C(5\tau)$ จากกราฟ
    \item เปรียบเทียบกับสมการ analytical
    \item อธิบายความคลาดเคลื่อนจาก numerical sampling หากมี
    \item จัดทำรายงานย่อ 2–4 หน้า
\end{enumerate}

\subsubsection{Mini Project: “One Problem, Three Tools”}

ให้ผู้เรียนเลือกหนึ่งหัวข้อจาก

\begin{itemize}
    \item เฟสเซอร์และอิมพีแดนซ์ตามความถี่
    \item เวกเตอร์และสนาม
    \item Fourier/FFT
    \item Laplace/RC/RLC
\end{itemize}

แล้วแก้ปัญหาด้วยอย่างน้อย 3 วิธี เช่น

\begin{enumerate}
    \item คำนวณด้วยมือ
    \item Python
    \item MATLAB/Octave หรือ GeoGebra
\end{enumerate}

จากนั้นเปรียบเทียบ

\begin{itemize}
    \item ความเร็ว
    \item ความโปร่งใสของขั้นตอน
    \item ความง่ายในการสร้างกราฟ
    \item ความง่ายในการตรวจสอบคำตอบ
    \item ข้อจำกัดของแต่ละเครื่องมือ
\end{itemize}

\par\noindent\rule{\textwidth}{0.4pt}\par
\section{8. ความเข้าใจผิดที่พบบ่อย}

\subsection{8.1 “โปรแกรมให้คำตอบแล้ว แปลว่าคำตอบถูก”}

ไม่ถูกต้อง โปรแกรมคำนวณตามคำสั่งที่ได้รับ หากกำหนดสมการผิด หน่วยผิด หรือ sampling ผิด โปรแกรมก็จะคำนวณสิ่งที่ผิดนั้นอย่างสม่ำเสมอ วิธีแก้คือใช้ \textbf{independent check} เช่น คำนวณหนึ่งจุดด้วยมือ ตรวจ limiting behavior หรือใช้เครื่องมืออีกชนิดตรวจซ้ำ

\subsection{8.2 “ผล symbolic ที่หน้าตาต่างกันคือคนละคำตอบ”}

นิพจน์

\[
\frac{1}{s+1/(RC)}
\]

และ

\[
\frac{RC}{RCs+1}
\]

เป็นสมการสมมูลกัน โปรแกรมแต่ละตัวอาจจัดรูปไม่เหมือนกัน ผู้เรียนควรจัดรูปหรือใช้การแทนค่าตรวจสอบก่อนสรุปว่าไม่ตรงกัน

\subsection{8.3 “FFT คือ Fourier series โดยตรง”}

FFT เป็นอัลกอริทึมคำนวณ DFT ของข้อมูลไม่ต่อเนื่องที่มีจำนวนตัวอย่างจำกัด ส่วน Fourier series เป็นการแทนฟังก์ชันเป็นคาบด้วยอนุกรมของฮาร์มอนิก แม้มีความเชื่อมโยงด้านความถี่ แต่เงื่อนไขและการตีความต่างกัน

\subsection{8.4 “FFT peak เท่ากับแอมพลิจูดสัญญาณเสมอ”}

ไม่เสมอ ต้องพิจารณา normalization, one-sided/two-sided spectrum, window, record length และตำแหน่งความถี่เทียบกับ FFT bins

\subsection{8.5 “Sampling rate สูงที่สุดเสมอดีที่สุด”}

sampling rate ที่สูงขึ้นช่วยครอบคลุมความถี่สูงขึ้น แต่เพิ่มจำนวนข้อมูล ภาระคำนวณ และอาจไม่ได้เพิ่ม frequency resolution หากระยะเวลาบันทึกไม่เปลี่ยน ความละเอียดความถี่ขึ้นกับ

\[
\Delta f=\frac{1}{T_{\text{record}}}
\]

จึงต้องออกแบบทั้ง $f_s$ และ $T_{\text{record}}$

\subsection{8.6 “หน่วยเป็นสิ่งที่โปรแกรมดูแลให้”}

โค้ดตัวเลขทั่วไปมักเห็นเพียงค่าตัวเลข เช่น \texttt{100e-6} โดยไม่รู้ว่านั่นคือฟารัด ถ้าผู้ใช้ใส่ \texttt{100} โปรแกรมก็อาจตีความเป็น 100 โดยไม่มีคำเตือน ดังนั้นควรตั้งชื่อตัวแปรและ comment ให้ชัดเจน

\subsection{8.7 “มุม phase จากโปรแกรมใช้ได้ทันทีทุกกรณี”}

ต้องตรวจหน่วยของมุมว่าคืนค่าเป็น radian หรือ degree และต้องระวัง phase wrapping เช่น $181^\circ$ กับ $-179^\circ$ อาจแทนทิศทางเกือบเดียวกัน

\subsection{8.8 “กราฟที่ smooth แปลว่าแบบจำลองแม่นยำ”}

ความเรียบของเส้นขึ้นกับจำนวนจุดและ interpolation แต่ความแม่นยำของแบบจำลองขึ้นกับสมมติฐานทางกายภาพ เช่น วงจรจริงมี ESR, parasitic inductance, tolerances และ bandwidth ของอุปกรณ์ ซึ่งอาจไม่อยู่ในแบบจำลองอุดมคติ

\par\noindent\rule{\textwidth}{0.4pt}\par

\section{9. การประยุกต์ใช้ในงานจริง}

\subsection{9.1 การตรวจสอบ power quality เบื้องต้น}

ข้อมูลแรงดันหรือกระแสจากเครื่องมือวัดสามารถนำเข้าโปรแกรมเพื่อ

\begin{enumerate}
    \item ตรวจ sampling information
    \item ตัดช่วงข้อมูลที่สนใจ
    \item ลบ DC offset หากมีเหตุผลรองรับ
    \item คำนวณ FFT
    \item อ่าน fundamental และ harmonic components
    \item คำนวณตัวชี้วัด เช่น THD ตามนิยามที่กำหนด
    \item สร้างกราฟเพื่อรายงาน
\end{enumerate}

หลักสำคัญคือข้อมูลจริงมี noise, leakage และความคลาดเคลื่อนของ sensor จึงไม่ควรตีความเหมือนสัญญาณสังเคราะห์ที่สมบูรณ์แบบ

\subsection{9.2 การออกแบบวงจรกรองเบื้องต้น}

สำหรับ RC low-pass

\[
f_c=\frac{1}{2\pi RC}
\]

โปรแกรมช่วย sweep ค่า $R$ และ $C$ แล้วสร้างกราฟ $|H(j\omega)|$ เพื่อดูผลต่อ cutoff frequency ได้ นักศึกษาสามารถกำหนดช่วงค่าจากอุปกรณ์มาตรฐาน แล้วเลือกค่าที่เข้าใกล้เป้าหมายที่สุด

\subsection{9.3 การตรวจสอบ transient ของวงจร}

วงจร RC, RL และ RLC สามารถเริ่มจาก ODE แล้ว

\begin{itemize}
    \item ใช้ Laplace หา analytical solution
    \item ใช้ SciPy/MATLAB คำนวณ step response
    \item เปรียบเทียบกราฟ
    \item วิเคราะห์ pole และ time constant
\end{itemize}

วิธีนี้เหมาะกับการตรวจการคำนวณและสร้าง intuition ก่อนใช้เครื่องมือจำลองวงจรเฉพาะทาง

\subsection{9.4 การสำรวจสนามไฟฟ้าและสนามแม่เหล็ก}

ถ้ามีศักย์ $V(x,y)$ หรือฟังก์ชันสนาม สามารถใช้ symbolic differentiation หา $\nabla V$, divergence หรือ curl และใช้กราฟ vector field/contour ช่วยให้เห็นทิศทางและความเข้มของสนาม

\subsection{9.5 การเตรียมข้อมูลสำหรับรายงานและงานวิจัย}

กระบวนการที่ดีควรแยก

\begin{itemize}
    \item raw data
    \item processed data
    \item script
    \item figures
    \item report
\end{itemize}

ออกจากกันอย่างมีระบบ และเก็บพารามิเตอร์สำคัญไว้ใน script หรือไฟล์ config เพื่อลดความผิดพลาดจากการแก้ตัวเลขหลายจุด

% [รูปที่ 9: โครงสร้างแฟ้มงานสำหรับ computational engineering report]
% ประเภทภาพ: Infographic
% ตำแหน่งที่แนะนำ: หัวข้อ 9.5
% วัตถุประสงค์การเรียนรู้: ส่งเสริม reproducibility และการจัดการไฟล์อย่างเป็นระบบ
% คำอธิบายภาพ: โครงสร้าง project folder มี data_raw/, data_processed/, src/, figures/, report/, README และลูกศร workflow
% องค์ประกอบที่ต้องมี: folder icons, script, CSV/data, plots, report
% ภาษาของป้ายกำกับ: อังกฤษเป็นชื่อโฟลเดอร์และไทยอธิบายสั้น ๆ
% สัดส่วนภาพ: 16:9
% Image Generation Prompt: Clean engineering infographic showing a reproducible computational project folder structure: data_raw, data_processed, src, figures, report, and README. Show workflow arrows from raw measurements to processing scripts to figures and final technical report. Minimal bilingual Thai-English annotations, white background, professional vector textbook style, 16:9.
% Alt Text: โครงสร้างโฟลเดอร์สำหรับเก็บข้อมูลดิบ โค้ด กราฟ และรายงานอย่างเป็นระบบ

\begin{figure}[htbp]
    \centering
    \includegraphics[width=0.8\textwidth]{images/figure-09.png}
    \caption{โครงสร้างแฟ้มงานสำหรับ computational engineering report}
\end{figure}

\par\noindent\rule{\textwidth}{0.4pt}\par

\section{10. สรุปท้ายบท}

บทนี้ใช้โปรแกรมคำนวณเป็นเครื่องมือบูรณาการคณิตศาสตร์วิศวกรรมไฟฟ้าจากบทก่อนหน้า โดยสาระสำคัญสรุปได้ดังนี้

\begin{enumerate}
    \item \textbf{เครื่องมือแต่ละชนิดมีจุดเด่นต่างกัน}
    เครื่องคิดเลขเหมาะกับการตรวจค่าเร็ว Python เหมาะกับการคำนวณเชิงตัวเลข/ข้อมูลและขยายได้ด้วยไลบรารี MATLAB/Octave เหมาะกับงานเมทริกซ์และงานวิศวกรรมแบบสคริปต์ ส่วน GeoGebra เหมาะกับการสำรวจกราฟและแนวคิดแบบโต้ตอบ
    \item \textbf{Numerical และ Symbolic เป็นคนละแนวทางที่เสริมกัน}
    numerical ให้ค่าประมาณจากข้อมูล ส่วน symbolic ให้รูปสมการ ผู้เรียนควรเลือกให้เหมาะกับปัญหา
    \item \textbf{จำนวนเชิงซ้อนเป็นฐานสำคัญของการคำนวณเฟสเซอร์และสเปกตรัม}
    โปรแกรมช่วยแปลง rectangular/polar และคำนวณอิมพีแดนซ์ได้รวดเร็ว แต่ต้องตรวจหน่วยมุมและความหมายทางกายภาพ
    \item \textbf{เวกเตอร์และ field operations สามารถคำนวณทั้งเชิงตัวเลขและสัญลักษณ์}
    การตรวจ orthogonality, gradient, divergence และ curl ช่วยเชื่อมคณิตศาสตร์กับสนามไฟฟ้าและสนามแม่เหล็ก
    \item \textbf{FFT ต้องใช้ความเข้าใจเรื่อง sampling}
    ความสัมพันธ์หลักคือ

    \[
    \Delta f=\frac{f_s}{N}=\frac{1}{T_{\text{record}}}
    \]

    และต้องพิจารณา Nyquist, leakage, normalization และ one-sided/two-sided spectrum
    \item \textbf{Laplace และ transfer function เชื่อมสมการวงจรกับการตอบสนองของระบบ}
    โปรแกรม symbolic ช่วยหาทรานส์ฟอร์ม ส่วน numerical tools ช่วยสร้าง response curve และตรวจ analytical solution
    \item \textbf{การตรวจสอบผลเป็นส่วนหนึ่งของการคำนวณ}
    ใช้การคำนวณด้วยมือ limiting behavior หน่วย และหลักฟิสิกส์เป็นตัวตรวจ
    \item \textbf{รายงานที่ดีต้องทำซ้ำได้}
    ต้องระบุสมการ พารามิเตอร์ หน่วย เครื่องมือ วิธีคำนวณ กราฟ การตรวจสอบ และข้อจำกัด
\end{enumerate}

\par\noindent\rule{\textwidth}{0.4pt}\par

\section{11. คำถามทบทวน}

\begin{enumerate}
    \item เพราะเหตุใดโปรแกรมคำนวณจึงไม่สามารถทดแทนความเข้าใจในแบบจำลองทางคณิตศาสตร์ได้?
    \item Numerical computation และ symbolic computation แตกต่างกันอย่างไร และยกตัวอย่างงานวิศวกรรมไฟฟ้าที่เหมาะกับแต่ละแบบ
    \item เหตุใดการใช้ \texttt{atan2} จึงเหมาะกับการหามุมของจำนวนเชิงซ้อนมากกว่าการใช้ $\tan^{-1}(b/a)$ เพียงอย่างเดียว?
    \item ถ้าใช้ Python คำนวณ phase แล้วได้ค่าเป็น radian แต่ผู้เรียนตีความเป็น degree ผลจะผิดในลักษณะใด?
    \item อธิบายความสัมพันธ์ระหว่าง sampling frequency, จำนวนตัวอย่าง, record duration และ frequency resolution
    \item เพราะเหตุใด FFT amplitude จึงต้องพิจารณา normalization ก่อนเทียบกับแอมพลิจูดจริงของสัญญาณ?
    \item อธิบาย spectral leakage และยกตัวอย่างสาเหตุที่ทำให้เกิด
    \item การหา Laplace transform ด้วยโปรแกรม symbolic ควรตรวจสอบคำตอบด้วยวิธีใดได้บ้าง?
    \item สำหรับวงจร RC low-pass ถ้ากราฟ step response จากโปรแกรมลดลงจาก 1 ไป 0 ควรสงสัยจุดใดของแบบจำลองหรือการกำหนด input/output?
    \item รายงานผลการคำนวณที่ “ทำซ้ำได้” ควรมีข้อมูลอะไรอย่างน้อย 5 รายการ?
\end{enumerate}

\par\noindent\rule{\textwidth}{0.4pt}\par

\section{12. แบบฝึกหัดท้ายบท}

\begin{quote}
\textbf{คำแนะนำ:} แบบฝึกหัดแบ่งเป็นระดับพื้นฐาน วิเคราะห์ และประยุกต์ใช้ตามลำดับ ให้ผู้เรียนแสดงสมการ หน่วย วิธีใช้โปรแกรม และการตรวจสอบผล ไม่ให้ส่งเฉพาะผลลัพธ์จากหน้าจอ
\end{quote}

\subsection{12.1 ระดับพื้นฐาน — 5 ข้อ}

\subsubsection{ข้อ 1 จำนวนเชิงซ้อน}

กำหนด

\[
Z=8-j6
\]

\begin{enumerate}
    \item หา $|Z|$ และ $\angle Z$ ด้วยมือ
    \item ตรวจคำตอบด้วย Python หรือเครื่องคิดเลขวิศวกรรม
    \item แปลงมุมเป็น degree
\end{enumerate}

\subsubsection{ข้อ 2 เฟสเซอร์กระแส}

กำหนด

\[
V=120\angle 20^\circ\ \text{V},\qquad Z=15\angle -30^\circ\ \Omega
\]

หาก

\[
I=\frac{V}{Z}
\]

จงหาขนาดและมุมของกระแส และตรวจด้วยโปรแกรม

\subsubsection{ข้อ 3 เวกเตอร์}

กำหนด

\[
\mathbf{A}=[1,2,-2],\qquad \mathbf{B}=[3,-1,4]
\]

ใช้ NumPy หรือ MATLAB/Octave หา

\begin{enumerate}
    \item $\mathbf{A}\cdot\mathbf{B}$
    \item $\mathbf{A}\times\mathbf{B}$
    \item ตรวจว่า cross product ตั้งฉากกับ $\mathbf{A}$ และ $\mathbf{B}$
\end{enumerate}

\subsubsection{ข้อ 4 สัญญาณไซน์}

สร้างสัญญาณ

\[
v(t)=100\sin(2\pi60t+30^\circ)
\]

ช่วง $0\le t\le0.05$ s ด้วย sampling frequency 5 kHz แล้วสร้างกราฟที่มีชื่อแกนและหน่วยครบถ้วน

\subsubsection{ข้อ 5 Laplace เชิงสัญลักษณ์}

ใช้ SymPy หรือ MATLAB Symbolic หา

\[
\mathcal{L}\{3e^{-2t}\}
\]

และตรวจคำตอบจากตารางลาปลาซ

\par\noindent\rule{\textwidth}{0.4pt}\par

\subsection{12.2 ระดับวิเคราะห์ — 5 ข้อ}

\subsubsection{ข้อ 6 อิมพีแดนซ์ตามความถี่}

วงจร RLC อนุกรมมี

\[
R=10\ \Omega,\qquad L=20\ \text{mH},\qquad C=220\ \mu\text{F}
\]

\begin{enumerate}
    \item คำนวณ $f_0$ ด้วยมือ
    \item ใช้โปรแกรม sweep $f=10$ ถึง 1000 Hz
    \item สร้างกราฟ $|Z(f)|$
    \item ระบุความถี่ที่ $|Z|$ ต่ำสุดจากข้อมูลเชิงตัวเลข
    \item อภิปรายความแตกต่างระหว่างค่าทฤษฎีและค่าที่อ่านจากกราฟ
\end{enumerate}

\subsubsection{ข้อ 7 Gradient และสนามไฟฟ้า}

กำหนดศักย์

\[
V(x,y)=5x^2+3xy-2y^2\quad \text{V}
\]

\begin{enumerate}
    \item ใช้การหาอนุพันธ์ด้วยมือหา $\nabla V$
    \item ใช้ SymPy ตรวจสอบ
    \item หา $\mathbf{E}=-\nabla V$ ที่จุด $(x,y)=(1,2)$
    \item อธิบายทิศของสนามเทียบกับการเพิ่มของศักย์
\end{enumerate}

\subsubsection{ข้อ 8 Frequency resolution}

ระบบเก็บข้อมูลด้วย

\[
f_s=4000\ \text{Hz}
\]

กรณี A บันทึก 0.25 s และกรณี B บันทึก 2 s

\begin{enumerate}
    \item หา $N$ ของแต่ละกรณี
    \item หา $\Delta f$
    \item ถ้าต้องแยกสัญญาณ 100 Hz และ 102 Hz กรณีใดเหมาะกว่า เพราะเหตุใด?
\end{enumerate}

\subsubsection{ข้อ 9 FFT และ normalization}

สร้าง

\[
x(t)=4\sin(2\pi100t)
\]

ที่ $f_s=2000$ Hz เป็นเวลา 1 s แล้ว

\begin{enumerate}
    \item คำนวณ \texttt{rfft}
    \item สร้าง single-sided amplitude spectrum
    \item อธิบายเหตุผลที่ต้องคูณ 2 สำหรับ bins ด้านบวกที่ไม่ใช่ DC และ Nyquist ตามเงื่อนไขที่เกี่ยวข้อง
    \item ตรวจว่า peak ใกล้ 100 Hz มีค่าแอมพลิจูดประมาณ 4
\end{enumerate}

\subsubsection{ข้อ 10 RC step response}

กำหนด

\[
R=4.7\ \text{k}\Omega,\qquad C=22\ \mu\text{F},\qquad V_s=10\ \text{V}
\]

\begin{enumerate}
    \item หา $\tau$
    \item เขียน $H(s)$
    \item หา $v_C(t)$ สำหรับ step input ด้วยมือ
    \item ใช้โปรแกรมสร้างกราฟช่วง $0$ ถึง $5\tau$
    \item เปรียบเทียบค่าที่ $t=\tau$ และ $t=5\tau$
\end{enumerate}

\par\noindent\rule{\textwidth}{0.4pt}\par

\subsection{12.3 ระดับประยุกต์ใช้ — 3 ข้อ}

\subsubsection{ข้อ 11 วิเคราะห์ฮาร์มอนิกจากข้อมูลสังเคราะห์}

สร้างสัญญาณแรงดัน

\[
v(t)=311\sin(2\pi50t)+20\sin(2\pi150t)+10\sin(2\pi250t)
\]

โดยใช้ $f_s=5000$ Hz และช่วงเวลา 1 s

\begin{enumerate}
    \item สร้าง time-domain plot ช่วง 0–0.06 s
    \item สร้าง single-sided spectrum 0–500 Hz
    \item ระบุ fundamental และ harmonics
    \item อภิปรายว่าเหตุใด 150 และ 250 Hz จึงสัมพันธ์กับฮาร์มอนิกลำดับที่ 3 และ 5 ของ 50 Hz
    \item หากลด record duration เหลือ 0.1 s จะเกิดผลต่อ frequency resolution อย่างไร?
\end{enumerate}

\subsubsection{ข้อ 12 เปรียบเทียบวิธีหาการตอบสนองวงจร}

สำหรับ RC low-pass ที่ $R=1.5\ \text{k}\Omega$, $C=68\ \mu\text{F}$ และ input step 5 V ให้หาการตอบสนองด้วย 3 วิธี

\begin{enumerate}
    \item analytical equation
    \item SymPy inverse Laplace
    \item SciPy หรือ MATLAB numerical step response
\end{enumerate}

สร้างตารางเปรียบเทียบค่าที่ $t=0$, $\tau$, $2\tau$, $5\tau$ และอภิปรายความแตกต่าง

\subsubsection{ข้อ 13 โครงงานย่อยเลือกเครื่องมือ}

เลือกปัญหาหนึ่งจาก จำนวนเชิงซ้อน/เฟสเซอร์, vector field, Fourier/FFT หรือ Laplace circuit แล้ว

\begin{enumerate}
    \item ระบุวัตถุประสงค์
    \item ระบุสมการและสมมติฐาน
    \item เลือกเครื่องมืออย่างน้อย 2 ชนิด พร้อมเหตุผล
    \item คำนวณหรือสร้างกราฟ
    \item ตรวจสอบผลอย่างน้อย 2 วิธี
    \item สรุปว่าเครื่องมือใดเหมาะกับงานมากกว่าและเพราะเหตุใด
\end{enumerate}

\par\noindent\rule{\textwidth}{0.4pt}\par

\section{12.4 แนวคำตอบและเฉลยแบบฝึกหัด}

\begin{quote}
ส่วนนี้เหมาะสำหรับผู้สอนหรือใช้หลังจากผู้เรียนส่งงานแล้ว
\end{quote}

\subsection{เฉลยข้อ 1}

ให้

\[
Z=8-j6
\]

ขนาด

\[
|Z|=\sqrt{8^2+(-6)^2}=10
\]

มุม

\[
\theta=\operatorname{atan2}(-6,8)\approx-36.87^\circ
\]

ดังนั้น

\[
Z=10\angle-36.87^\circ
\]

ตัวอย่าง Python

\begin{verbatim}
import cmath, math
z = 8 - 6j
print(abs(z))
print(math.degrees(cmath.phase(z)))
\end{verbatim}

\textbf{ตรวจความสมเหตุสมผล:} $Z$ อยู่ควอดแรนต์ที่ 4 เพราะส่วนจริงบวกและส่วนจินตภาพลบ จึงต้องได้มุมลบ

\subsection{เฉลยข้อ 2}

\[
I=\frac{120\angle20^\circ}{15\angle-30^\circ}
\]

หารขนาดและลบมุม

\[
|I|=\frac{120}{15}=8\ \text{A}
\]

\[
\angle I=20^\circ-(-30^\circ)=50^\circ
\]

จึงได้

\[
I=8\angle50^\circ\ \text{A}
\]

\textbf{ข้อผิดพลาดที่พบบ่อย:} บวกมุมของตัวหารแทนการลบมุม หรือใช้ degree โดยไม่แปลงเป็น radian ในฟังก์ชัน trigonometric ของโปรแกรม

\subsection{เฉลยข้อ 3}

ให้

\[
\mathbf{A}=[1,2,-2],\quad \mathbf{B}=[3,-1,4]
\]

ดอตโปรดักต์

\[
\mathbf{A}\cdot\mathbf{B}
=1(3)+2(-1)+(-2)(4)=3-2-8=-7
\]

ครอสโปรดักต์

\[
\mathbf{A}\times\mathbf{B}
=
\begin{vmatrix}
\hat x&\hat y&\hat z\\
1&2&-2\\
3&-1&4
\end{vmatrix}
\]

ได้

\[
\mathbf{A}\times\mathbf{B}=[6,-10,-7]
\]

ตรวจ

\[
[6,-10,-7]\cdot[1,2,-2]=6-20+14=0
\]

และ

\[
[6,-10,-7]\cdot[3,-1,4]=18+10-28=0
\]

ดังนั้น cross product ตั้งฉากกับทั้งสองเวกเตอร์

\subsection{เฉลยข้อ 4}

แปลงเฟส 30° เป็นเรเดียน

\[
30^\circ=\frac{\pi}{6}\ \text{rad}
\]

ตัวอย่าง Python

\begin{verbatim}
import numpy as np
import matplotlib.pyplot as plt

fs = 5000
T = 0.05
t = np.arange(0, T, 1/fs)
v = 100*np.sin(2*np.pi*60*t + np.deg2rad(30))

plt.plot(t,v)
plt.xlabel("Time (s)")
plt.ylabel("Voltage (V)")
plt.title("60-Hz Sinusoid, phase = 30 deg")
plt.grid(True)
plt.show()
\end{verbatim}

จำนวนตัวอย่างโดยประมาณ

\[
N=f_sT=5000(0.05)=250
\]

\subsection{เฉลยข้อ 5}

จากตาราง

\[
\mathcal{L}\{e^{-at}\}=\frac{1}{s+a}
\]

ดังนั้น

\[
\mathcal{L}\{3e^{-2t}\}=\frac{3}{s+2}
\]

SymPy

\begin{verbatim}
import sympy as sp

t, s = sp.symbols('t s', positive=True)
F = sp.laplace_transform(3*sp.exp(-2*t), t, s, noconds=True)
print(F)
\end{verbatim}

\subsection{เฉลยข้อ 6}

กำหนด

\[
R=10\ \Omega,\quad L=0.02\ \text{H},\quad C=220\times10^{-6}\ \text{F}
\]

เรโซแนนซ์

\[
f_0=\frac{1}{2\pi\sqrt{LC}}
\]

แทนค่า

\[
LC=(0.02)(220\times10^{-6})=4.4\times10^{-6}
\]

\[
\sqrt{LC}\approx2.0976\times10^{-3}
\]

จึงได้

\[
f_0\approx75.86\ \text{Hz}
\]

ตัวอย่างโค้ด

\begin{verbatim}
import numpy as np
import matplotlib.pyplot as plt

R = 10.0
L = 20e-3
C = 220e-6
f = np.linspace(10, 1000, 5000)
w = 2*np.pi*f
Z = R + 1j*(w*L - 1/(w*C))

idx = np.argmin(np.abs(Z))
print(f[idx], np.abs(Z[idx]))

plt.plot(f, np.abs(Z))
plt.xlabel("Frequency (Hz)")
plt.ylabel("|Z| (ohm)")
plt.grid(True)
plt.show()
\end{verbatim}

ค่าที่อ่านจากข้อมูลอาจไม่ตรง 75.86 Hz พอดี เพราะ grid ความถี่เป็นค่าจำกัด หากเพิ่มจำนวนจุดหรือคำนวณที่ $f_0$ โดยตรงจะใกล้ค่าทฤษฎีมากขึ้น

\subsection{เฉลยข้อ 7}

\[
V(x,y)=5x^2+3xy-2y^2
\]

ดังนั้น

\[
\frac{\partial V}{\partial x}=10x+3y
\]

\[
\frac{\partial V}{\partial y}=3x-4y
\]

จึงได้

\[
\nabla V=(10x+3y)\hat x+(3x-4y)\hat y
\]

สนามไฟฟ้า

\[
\mathbf{E}=-(10x+3y)\hat x-(3x-4y)\hat y
\]

ที่ $(1,2)$

\[
\mathbf{E}=-(10+6)\hat x-(3-8)\hat y
\]

\[
\boxed{\mathbf{E}=-16\hat x+5\hat y\ \text{V/m}}
\]

ทิศของ $\mathbf{E}$ ตรงข้ามกับทิศที่ศักย์เพิ่มเร็วที่สุด

\subsection{เฉลยข้อ 8}

กรณี A

\[
N_A=f_sT=4000(0.25)=1000
\]

\[
\Delta f_A=\frac{1}{0.25}=4\ \text{Hz}
\]

กรณี B

\[
N_B=4000(2)=8000
\]

\[
\Delta f_B=\frac{1}{2}=0.5\ \text{Hz}
\]

ต้องการแยก 100 และ 102 Hz ซึ่งห่างกัน 2 Hz กรณี B เหมาะกว่าเพราะความละเอียด 0.5 Hz ขณะที่กรณี A มี bin spacing 4 Hz จึงแยกสององค์ประกอบนี้ได้ยากกว่า

\subsection{เฉลยข้อ 9}

สำหรับ $T=1$ s

\[
N=2000
\]

และ

\[
\Delta f=1\ \text{Hz}
\]

100 Hz จึงตรงกับ bin ได้ดี ตัวอย่าง

\begin{verbatim}
import numpy as np

fs = 2000.0
T = 1.0
t = np.arange(0, T, 1/fs)
x = 4*np.sin(2*np.pi*100*t)

N = len(x)
X = np.fft.rfft(x)
f = np.fft.rfftfreq(N, 1/fs)
A = np.abs(X)/N
A[1:-1] *= 2

k = np.argmax(A)
print(f[k], A[k])
\end{verbatim}

แนวคิด normalization: DFT ของสัญญาณจริงกระจายองค์ประกอบไซน์/โคไซน์ไปยังความถี่บวกและลบ ใน single-sided spectrum จึงรวมพลังของคู่ความถี่โดยคูณสองใน bins ที่เหมาะสม ส่วน DC และ Nyquist ต้องจัดการตามโครงสร้างข้อมูลและเงื่อนไขของ $N$

\subsection{เฉลยข้อ 10}

\[
R=4700\ \Omega,\quad C=22\times10^{-6}\ \text{F}
\]

\[
\tau=RC=4700(22\times10^{-6})=0.1034\ \text{s}
\]

ฟังก์ชันถ่ายโอน

\[
H(s)=\frac{1}{0.1034s+1}
\]

สำหรับ step 10 V

\[
v_C(t)=10\left(1-e^{-t/0.1034}\right)
\]

ที่ $t=\tau$

\[
v_C(\tau)=10(1-e^{-1})\approx6.321\ \text{V}
\]

ที่ $t=5\tau$

\[
v_C(5\tau)=10(1-e^{-5})\approx9.933\ \text{V}
\]

กราฟจากโปรแกรมควรผ่านค่าดังกล่าวโดยประมาณ

\subsection{เฉลยข้อ 11}

สัญญาณประกอบด้วย

\begin{itemize}
    \item 50 Hz แอมพลิจูด 311 V → fundamental
    \item 150 Hz แอมพลิจูด 20 V → $3\times50$ Hz = 3rd harmonic
    \item 250 Hz แอมพลิจูด 10 V → $5\times50$ Hz = 5th harmonic
\end{itemize}

เมื่อ $T=1$ s

\[
\Delta f=1\ \text{Hz}
\]

ถ้าลดเป็น $T=0.1$ s

\[
\Delta f=10\ \text{Hz}
\]

จึงมี frequency resolution หยาบขึ้น แม้ sampling rate เดิมยังคงครอบคลุมความถี่เหล่านี้ได้

ตัวอย่างโค้ด

\begin{verbatim}
import numpy as np
import matplotlib.pyplot as plt

fs = 5000.0
T = 1.0
t = np.arange(0, T, 1/fs)
v = (311*np.sin(2*np.pi*50*t)
     + 20*np.sin(2*np.pi*150*t)
     + 10*np.sin(2*np.pi*250*t))

N = len(v)
V = np.fft.rfft(v)
f = np.fft.rfftfreq(N, 1/fs)
A = np.abs(V)/N
A[1:-1] *= 2

plt.plot(t, v)
plt.xlim(0, 0.06)
plt.xlabel("Time (s)")
plt.ylabel("Voltage (V)")
plt.grid(True)
plt.show()

plt.stem(f, A)
plt.xlim(0, 500)
plt.xlabel("Frequency (Hz)")
plt.ylabel("Amplitude (V)")
plt.grid(True)
plt.show()
\end{verbatim}

\subsection{เฉลยข้อ 12}

\[
R=1500\ \Omega,\quad C=68\times10^{-6}\ \text{F}
\]

\[
\tau=RC=1500(68\times10^{-6})=0.102\ \text{s}
\]

คำตอบ analytical

\[
v_C(t)=5(1-e^{-t/0.102})
\]

ค่ามาตรฐาน

\begin{table}[htbp]
    \centering
    \begin{tabular}{l l r}
        \hline
        \textbf{เวลา} & \textbf{นิพจน์} & \textbf{$v_C$ โดยประมาณ} \\
        \hline
        $0$ & $5(1-e^0)$ & 0 V \\
        $\tau$ & $5(1-e^{-1})$ & 3.161 V \\
        $2\tau$ & $5(1-e^{-2})$ & 4.323 V \\
        $5\tau$ & $5(1-e^{-5})$ & 4.966 V \\
        \hline
    \end{tabular}
\end{table}

SymPy ควรให้รูปสมการสมมูล และ SciPy/MATLAB step response ควรให้ค่าที่ใกล้เคียงเมื่อใช้ time grid ละเอียดพอ

\textbf{การอภิปราย:} ความต่างเล็กน้อยในค่าที่อ่านจากกราฟมักเกิดจากการที่ sample time ไม่ตรงกับ $\tau$, $2\tau$, $5\tau$ พอดี หรือเกิดจากการ interpolation/precision ไม่ใช่เพราะทฤษฎีต่างกัน

\subsection{แนวคำตอบข้อ 13}

เป็นโจทย์เปิด ไม่มีคำตอบตัวเลขตายตัว แต่รายงานที่ดีควรมี

\begin{itemize}
    \item ปัญหาที่อยู่ในขอบเขตบท
    \item สมการถูกต้อง
    \item หน่วยครบ
    \item เหตุผลการเลือกเครื่องมือ
    \item ผลคำนวณหรือกราฟที่ตรวจสอบได้
    \item การตรวจอย่างน้อยสองวิธี เช่น มือ + software, analytical + numerical, limiting behavior + second tool
    \item ข้อจำกัดและข้อสรุปที่สัมพันธ์กับวัตถุประสงค์
\end{itemize}

\par\noindent\rule{\textwidth}{0.4pt}\par

\section{13. ข้อเสนอแนะสำหรับผู้สอน}

\subsection{13.1 ลำดับการสอน 2 สัปดาห์ที่แนะนำ}

\subsubsection{สัปดาห์ที่ 14: เครื่องมือและการคำนวณพื้นฐาน}

\textbf{ช่วงที่ 1:} บทบาทและหลัก verification \\
เริ่มจากโจทย์ RLC ตัวอย่างเดียวกัน ให้ผู้เรียนคำนวณด้วยมือก่อน แล้วใช้เครื่องคิดเลข/Python เพื่อยืนยัน

\textbf{ช่วงที่ 2:} NumPy, Matplotlib และจำนวนเชิงซ้อน \\
เปลี่ยนจากค่าความถี่เดียวเป็น sweep ความถี่ เพื่อให้เห็นเหตุผลว่าทำไมต้องใช้โปรแกรม

\textbf{ช่วงที่ 3:} SymPy และ vector \\
ใช้ gradient ตัวอย่างสั้น ๆ และตรวจด้วยมือหนึ่งองค์ประกอบ

\subsubsection{สัปดาห์ที่ 15: Fourier, Laplace และรายงาน}

\textbf{ช่วงที่ 1:} sampling + FFT \\
ให้สร้างสัญญาณสองความถี่ที่รู้คำตอบก่อน เพื่อให้ผู้เรียนเห็น peak ที่ควรเกิด

\textbf{ช่วงที่ 2:} RC symbolic vs numerical \\
ใช้โจทย์ RC เดียวกันใน SymPy และ SciPy/MATLAB แล้ววางกราฟทับกัน

\textbf{ช่วงที่ 3:} mini report \\
ให้นักศึกษาจัดทำรายงานย่อที่มีสมการ กราฟ การตรวจสอบ และข้อจำกัด

\subsection{13.2 จุดที่ผู้เรียนมักมีปัญหา}

\begin{enumerate}
    \item radian กับ degree
    \item prefix ของหน่วย เช่น m, $\mu$, k
    \item การสร้าง time vector ให้จำนวนจุดถูกต้อง
    \item การตีความ FFT bins
    \item normalization ของ spectrum
    \item symbolic answer ที่อยู่คนละรูป
    \item การกำหนด numerator/denominator ของ transfer function ผิดลำดับ
    \item การลืม zero initial condition เมื่อใช้ transfer function
    \item การส่ง screenshot แทนการอธิบายวิธีคำนวณ
    \item การรายงานกราฟโดยไม่มีหน่วย
\end{enumerate}

\subsection{13.3 วิธีประเมินระหว่างเรียน}

ใช้คำถามสั้นก่อนรันโปรแกรม เช่น

\begin{itemize}
    \item “ก่อนคำนวณ คิดว่ามุม $Z$ เป็นบวกหรือลบ?”
    \item “ก่อน FFT ควรเห็น peak ที่ความถี่ใด?”
    \item “ถ้าเพิ่ม $R$ ใน RLC การสั่นควรเพิ่มหรือลด?”
    \item “ถ้าเพิ่ม record duration แต่ $f_s$ คงเดิม $\Delta f$ เปลี่ยนอย่างไร?”
\end{itemize}

คำถามแบบ prediction ช่วยแยกผู้เรียนที่เข้าใจระบบออกจากผู้เรียนที่เพียงรันโค้ดได้

\subsection{13.4 Rubric งานรายงานย่อ}

\begin{table}[htbp]
    \centering
    \begin{tabular}{l r l}
        \hline
        \textbf{เกณฑ์} & \textbf{น้ำหนัก} & \textbf{ตัวบ่งชี้} \\
        \hline
        แบบจำลองและสมการ & 20\% & สมการถูกต้อง ระบุสมมติฐาน \\
        การใช้เครื่องมือ & 20\% & เลือกคำสั่งเหมาะสม โค้ดเป็นระเบียบ \\
        ผลลัพธ์และกราฟ & 20\% & หน่วย แกน legend และค่าจุดสำคัญครบ \\
        การตรวจสอบ & 20\% & มีมือ/ทฤษฎี/เหตุผลทางกายภาพตรวจซ้ำ \\
        การอภิปรายและข้อจำกัด & 10\% & แปลผล ไม่ใช่เพียงบรรยายกราฟ \\
        การสื่อสารและ reproducibility & 10\% & ไฟล์ครบ ทำซ้ำได้ อ้างอิงแหล่งข้อมูล \\
        \hline
    \end{tabular}
\end{table}

\par\noindent\rule{\textwidth}{0.4pt}\par

\section{14. สื่อและแหล่งเรียนรู้เพิ่มเติม}

\subsection{14.1 สื่อแนะนำ}

\subsubsection{Python Software Foundation}

\begin{itemize}
    \item \textbf{หัวข้อที่ควรค้น:} Python complex numbers, \texttt{cmath}
    \item \textbf{คำค้นแนะนำ:} \texttt{Python cmath complex polar phase}
    \item \textbf{จุดประสงค์:} ตรวจ syntax และฟังก์ชันสำหรับจำนวนเชิงซ้อน
    \item \textbf{ช่วงเวลาที่เหมาะสม:} ก่อนกิจกรรมเฟสเซอร์
\end{itemize}

\subsubsection{NumPy Documentation}

\begin{itemize}
    \item \textbf{หัวข้อที่ควรค้น:} Discrete Fourier Transform, \texttt{numpy.fft}, \texttt{rfft}, \texttt{rfftfreq}
    \item \textbf{คำค้นแนะนำ:} \texttt{NumPy FFT rfft frequency bins}
    \item \textbf{จุดประสงค์:} ตรวจความหมายของ FFT functions และการสร้างแกนความถี่
    \item \textbf{ช่วงเวลาที่เหมาะสม:} ก่อนกิจกรรม Fourier/FFT
\end{itemize}

\subsubsection{SciPy Documentation}

\begin{itemize}
    \item \textbf{หัวข้อที่ควรค้น:} \texttt{scipy.signal}, \texttt{TransferFunction}, \texttt{step}, \texttt{freqresp}
    \item \textbf{คำค้นแนะนำ:} \texttt{SciPy signal continuous transfer function step response}
    \item \textbf{จุดประสงค์:} เรียนรู้การแทนระบบ LTI และการตอบสนอง
    \item \textbf{ช่วงเวลาที่เหมาะสม:} ก่อนกิจกรรม RC/RLC
\end{itemize}

\subsubsection{Matplotlib Documentation}

\begin{itemize}
    \item \textbf{หัวข้อที่ควรค้น:} Pyplot tutorial, line plot, stem plot, labeling
    \item \textbf{คำค้นแนะนำ:} \texttt{Matplotlib pyplot plot stem xlabel ylabel}
    \item \textbf{จุดประสงค์:} สร้างกราฟที่มีองค์ประกอบครบสำหรับรายงาน
    \item \textbf{ช่วงเวลาที่เหมาะสม:} สัปดาห์แรกของบท
\end{itemize}

\subsubsection{SymPy Documentation}

\begin{itemize}
    \item \textbf{หัวข้อที่ควรค้น:} symbolic mathematics, Laplace transform, inverse Laplace transform
    \item \textbf{คำค้นแนะนำ:} \texttt{SymPy laplace\_transform inverse\_laplace\_transform}
    \item \textbf{จุดประสงค์:} ตรวจการคำนวณเชิงสัญลักษณ์และเปรียบเทียบกับการทำมือ
    \item \textbf{ช่วงเวลาที่เหมาะสม:} ก่อนกิจกรรม Laplace
\end{itemize}

\subsubsection{MathWorks Documentation}

\begin{itemize}
    \item \textbf{หัวข้อที่ควรค้น:} MATLAB \texttt{fft}, symbolic \texttt{laplace}, \texttt{ilaplace}
    \item \textbf{คำค้นแนะนำ:} \texttt{MATLAB fft laplace ilaplace}
    \item \textbf{จุดประสงค์:} เรียนรู้คำสั่งมาตรฐานในสภาพแวดล้อม MATLAB
    \item \textbf{ช่วงเวลาที่เหมาะสม:} เมื่อสถาบันหรือผู้เรียนใช้ MATLAB
\end{itemize}

\subsubsection{GNU Octave Documentation}

\begin{itemize}
    \item \textbf{หัวข้อที่ควรค้น:} vectors, matrices, plotting, FFT
    \item \textbf{คำค้นแนะนำ:} \texttt{GNU Octave manual fft plot complex}
    \item \textbf{จุดประสงค์:} ใช้เป็นทางเลือกโอเพนซอร์สสำหรับการคำนวณแบบ matrix-oriented
    \item \textbf{ช่วงเวลาที่เหมาะสม:} ใช้คู่กับตัวอย่าง MATLAB
\end{itemize}

\subsubsection{GeoGebra Team}

\begin{itemize}
    \item \textbf{หัวข้อที่ควรค้น:} Graphing Calculator, CAS Calculator, sliders
    \item \textbf{คำค้นแนะนำ:} \texttt{GeoGebra Graphing Calculator slider CAS}
    \item \textbf{จุดประสงค์:} สำรวจผลของพารามิเตอร์ต่อกราฟแบบ interactive
    \item \textbf{ช่วงเวลาที่เหมาะสม:} ก่อนหรือระหว่างกิจกรรม Fourier Explorer และ RC response
\end{itemize}

\par\noindent\rule{\textwidth}{0.4pt}\par

\section{15. เอกสารอ้างอิงท้ายบท}

\begin{quote}
รูปแบบอ้างอิง: APA 7 โดยรายการเว็บไซต์ใช้เอกสารทางการของผู้พัฒนาซอฟต์แวร์ และควรตรวจสอบวันที่เข้าถึง/รายละเอียดเวอร์ชันอีกครั้งก่อนจัดพิมพ์ฉบับเผยแพร่
\end{quote}

\subsection{15.1 เอกสารหลักที่ผู้ใช้แนบ}

โครงสร้างเอกสารประกอบการสอน รายวิชา 5571102 คณิตศาสตร์วิศวกรรมไฟฟ้า. (2569). \textit{โครงสร้างเอกสารประกอบการสอนและแผนการจัดการเรียนรู้รายสัปดาห์}.

Prompt กลางสำหรับออกแบบเอกสารประกอบการสอนระดับอุดมศึกษา. (ม.ป.ป.). \textit{แนวทางออกแบบเอกสารประกอบการสอนแบบ Topic-Based Design Mode}.

Module A. (ม.ป.ป.). \textit{ภาพประกอบและ Prompt สร้างภาพ}.

Module B. (ม.ป.ป.). \textit{สมการและตัวอย่างคำนวณ}.

Module E. (ม.ป.ป.). \textit{แบบฝึกหัดและเฉลย}.

Module F. (ม.ป.ป.). \textit{OBE และการประเมินผล}.

Module H. (ม.ป.ป.). \textit{สื่อ วิดีโอ และเอกสารอ้างอิง}.

\subsection{15.2 หนังสือและตำรา}

Kreyszig, E. (2011). \textit{Advanced engineering mathematics} (10th ed.). Wiley.

Oppenheim, A. V., Willsky, A. S., \& Nawab, S. H. (1997). \textit{Signals and systems} (2nd ed.). Prentice Hall.

Alexander, C. K., \& Sadiku, M. N. O. (2017). \textit{Fundamentals of electric circuits} (6th ed.). McGraw-Hill Education.

Chapra, S. C., \& Canale, R. P. (2015). \textit{Numerical methods for engineers} (7th ed.). McGraw-Hill Education.

\subsection{15.3 คู่มือและเอกสารทางเทคนิคอย่างเป็นทางการ}

Python Software Foundation. (2026). \textit{Python documentation: cmath—Mathematical functions for complex numbers}. เอกสารทางการออนไลน์ ตรวจสอบเมื่อ 19 สิงหาคม 2026.

NumPy Developers. (2026). \textit{NumPy documentation: Discrete Fourier Transform and numpy.fft}. เอกสารทางการออนไลน์ ตรวจสอบเมื่อ 19 สิงหาคม 2026.

SciPy Community. (2026). \textit{SciPy documentation: Signal processing (scipy.signal), continuous-time systems, step and frequency response}. เอกสารทางการออนไลน์ ตรวจสอบเมื่อ 19 สิงหาคม 2026.

Matplotlib Development Team. (2026). \textit{Matplotlib documentation: Pyplot and plotting APIs}. เอกสารทางการออนไลน์ ตรวจสอบเมื่อ 19 สิงหาคม 2026.

SymPy Development Team. (2026). \textit{SymPy documentation: Symbolic mathematics and Laplace transforms}. เอกสารทางการออนไลน์ ตรวจสอบเมื่อ 19 สิงหาคม 2026.

MathWorks. (2026). \textit{MATLAB documentation: FFT, Laplace transform, and inverse Laplace transform}. เอกสารทางการออนไลน์ ตรวจสอบเมื่อ 19 สิงหาคม 2026.

GNU Project. (2026). \textit{GNU Octave documentation}. เอกสารทางการออนไลน์ ตรวจสอบเมื่อ 19 สิงหาคม 2026.

GeoGebra Team. (2026). \textit{GeoGebra Graphing Calculator and Calculator Suite documentation}. เอกสารทางการออนไลน์ ตรวจสอบเมื่อ 19 สิงหาคม 2026.

\par\noindent\rule{\textwidth}{0.4pt}\par

\section{ภาคผนวก A: Quick Reference สำหรับ Python}

\subsection{A.1 จำนวนเชิงซ้อน}

\begin{verbatim}
import cmath
import math

z = 3 + 4j
mag = abs(z)
phase_rad = cmath.phase(z)
phase_deg = math.degrees(phase_rad)
r, phi = cmath.polar(z)
z2 = cmath.rect(r, phi)
\end{verbatim}

\subsection{A.2 NumPy array และสัญญาณ}

\begin{verbatim}
import numpy as np

fs = 5000
T = 1.0
t = np.arange(0, T, 1/fs)
x = np.sin(2*np.pi*50*t)
\end{verbatim}

\subsection{A.3 FFT แบบสัญญาณจริง}

\begin{verbatim}
N = len(x)
X = np.fft.rfft(x)
f = np.fft.rfftfreq(N, d=1/fs)
A = np.abs(X)/N
A[1:-1] *= 2
\end{verbatim}

\begin{quote}
\textbf{หมายเหตุ:} สูตร normalization ต้องพิจารณาว่า $N$ เป็นเลขคู่/คี่ และมี Nyquist bin หรือไม่ ตัวอย่างนี้ใช้เพื่อการเรียนรู้ในกรณีพื้นฐาน
\end{quote}

\subsection{A.4 กราฟ}

\begin{verbatim}
import matplotlib.pyplot as plt

plt.plot(t, x)
plt.xlabel("Time (s)")
plt.ylabel("Amplitude")
plt.grid(True)
plt.show()
\end{verbatim}

\subsection{A.5 Symbolic Laplace}

\begin{verbatim}
import sympy as sp

t, s = sp.symbols('t s', positive=True)
f = sp.exp(-2*t)
F = sp.laplace_transform(f, t, s, noconds=True)
f_back = sp.inverse_laplace_transform(F, s, t)
\end{verbatim}

\subsection{A.6 SciPy transfer function}

\begin{verbatim}
from scipy import signal

num = [1]
den = [0.1, 1]
sys = signal.TransferFunction(num, den)
t, y = signal.step(sys)
\end{verbatim}

\par\noindent\rule{\textwidth}{0.4pt}\par

\section{ภาคผนวก B: Quick Reference สำหรับ MATLAB / GNU Octave}

\subsection{B.1 Complex}

\begin{verbatim}
z = 3 + 4i;
mag = abs(z);
phase_rad = angle(z);
phase_deg = phase_rad*180/pi;
\end{verbatim}

\subsection{B.2 สัญญาณและกราฟ}

\begin{verbatim}
fs = 5000;
T = 1;
t = 0:1/fs:T-1/fs;
x = sin(2*pi*50*t);
plot(t,x);
xlabel('Time (s)');
ylabel('Amplitude');
grid on;
\end{verbatim}

\subsection{B.3 FFT}

\begin{verbatim}
N = length(x);
X = fft(x);
f = (0:N-1)*(fs/N);
\end{verbatim}

\subsection{B.4 MATLAB Symbolic Laplace}

\begin{verbatim}
syms t s
f = exp(-2*t);
F = laplace(f,t,s);
f2 = ilaplace(F,s,t);
\end{verbatim}

\begin{quote}
สำหรับ GNU Octave ความพร้อมของ symbolic operations ขึ้นกับ packages ที่ติดตั้ง จึงควรตรวจคู่มือและ package documentation ของระบบที่ใช้จริง
\end{quote}

\par\noindent\rule{\textwidth}{0.4pt}\par

\section{ภาคผนวก C: Checklist ก่อนส่งงานคำนวณ}

\begin{itemize}
    \item {}[ ] ระบุโจทย์และวัตถุประสงค์ชัดเจน
    \item {}[ ] สมการสอดคล้องกับระบบจริง
    \item {}[ ] ระบุสมมติฐาน
    \item {}[ ] ใช้หน่วย SI อย่างสม่ำเสมอ
    \item {}[ ] แปลง prefix m, µ, k, M ถูกต้อง
    \item {}[ ] ตรวจ degree/radian
    \item {}[ ] ระบุ sampling frequency และ record duration เมื่อมีสัญญาณ
    \item {}[ ] ตรวจ frequency resolution เมื่อใช้ FFT
    \item {}[ ] กำหนด normalization ของ spectrum ชัดเจน
    \item {}[ ] กราฟมีชื่อแกนและหน่วย
    \item {}[ ] ตรวจคำตอบอย่างน้อยหนึ่งวิธีที่เป็นอิสระจากโปรแกรมหลัก
    \item {}[ ] อธิบายความหมายทางวิศวกรรมของผลลัพธ์
    \item {}[ ] ระบุข้อจำกัดของแบบจำลอง
    \item {}[ ] แนบ script/notebook หรือไฟล์ต้นฉบับที่จำเป็น
    \item {}[ ] รายงานสามารถทำซ้ำได้จากข้อมูลที่ให้
\end{itemize}

\par\noindent\rule{\textwidth}{0.4pt}\par

\textbf{จบบทที่ 8 โปรแกรมคำนวณทางคณิตศาสตร์สำหรับวิศวกรรมไฟฟ้า}
